\subsection{\texorpdfstring{类型为 \(\mathbf{B}_l\) 的代数 ( \(l \geq 1\) )}{类型为 \textbackslash mathbf\{B\}\_l 的代数 ( l \textbackslash geq 1 )}}

\begin{enumerate}
\def\labelenumi{(\Roman{enumi})}
\tightlist
\item
  设 V 为有限维向量空间，\(\Psi\) 为 V 上的非退化对称双线性型。所有满足 \(\Psi(xv,v') + \Psi(v, xv') = 0\)（对所有 \(v, v' \in V\)）的 V 的自同态 x 构成 \(\mathfrak{sl}(V)\) 的李子代数；当 dim \(V \neq 2\) 时，它是半单的（第 I 章，§6，第7节，命题9）。我们将其记为 \(\mathfrak{o}(\Psi)\)，称其为与 \(\Psi\) 相关的正交李代数。
\end{enumerate}

设 V 的维数为奇数 \(2l + 1 \geq 3\)，且 \(\Psi\) 的指标为最大指标 l。记 Q 为与 \(\Psi\) 相关的二次型。有 \(\mathrm{Q}(x)=\frac{1}{2}\varPsi(x,x)\) 对 \(x\in V\) 成立。由《代数》第 IX 章，§4，第2节，V 可写成两个极大全迷向子空间 F 和 \(F'\) 与 \(F+F'\) 的正交补 G 的直和；G 非迷向且为一维。将 \(\varPsi\) 乘以一个非零常数后，不妨设存在 \(e_{0}\in G\)，使得 \(\varPsi(e_{0},e_{0})=-2\)。另一方面，F 和 \(F'\) 通过 \(\varPsi\) 处于对偶关系；令 \((e_{i})_{1\leq i\leq l}\) 为 F 的一个基，令 \((e_{-i})_{1\leq i\leq l}\) 为 \(F'\) 的对偶基。于是

\[
(e _ {1}, \ldots , e _ {l}, e _ {0}, e _ {- l}, \ldots , e _ {- 1})
\]

是 V 的一个基；有

\[
\mathrm{Q} \left(\sum x _ {i} e _ {i}\right) = - x _ {0} ^ {2} + \sum_ {i = 1} ^ {i = l} x _ {i} x _ {- i}
\]

并且相对于此基，\(\varPsi\) 的矩阵是阶为 \(2l + 1\) 的方阵

\[
S = \left( \begin{array}{c c c} 0 & 0 & s \\ 0 & - 2 & 0 \\ s & 0 & 0 \end{array} \right), \qquad s = \left( \begin{array}{c c c c c} 0 & 0 & \dots & 0 & 1 \\ 0 & 0 & \dots & 1 & 0 \\ \vdots & \vdots & \ddots & \vdots & \vdots \\ 0 & 1 & \dots & 0 & 0 \\ 1 & 0 & \dots & 0 & 0 \end{array} \right),
\]

其中 s 是阶为 l 的方阵，除第二对角线 \(^{6}\) 上的元素等于 1 外，其余元素全为 0。具有上述性质的 V 的基称为 Witt 基。于是，代数 \(\mathfrak{g} = \mathfrak{o}(\varPsi)\) 可与方阵代数 \(\mathfrak{o}_{S}(2l + 1, k)\) 识别，后者由阶为 \(2l + 1\) 的方阵 a 组成，且满足 \(a = -S^{-1}aS\)（《代数》第 IX 章，§1，第10节，公式（50））。直接计算表明，g 是如下形式的矩阵的集合：

\[
\left( \begin{array}{c c c} A & 2 s ^ {t} x & B \\ y & 0 & x \\ C & 2 s ^ {t} y & D \end{array} \right)\tag{2}
\]

其中 x 和 y 是具有 1 行 l 列的矩阵，A、B、C、D 是阶为 l 的方阵，并满足 \(B = -s^{t}Bs\)、\(C = -s^{t}Cs\) 和 \(D = -s^{t}As\)。由于从 \(A \mapsto s^{t}As\) 到其自身的映射 \(\mathbf{M}_{l}(k)\) 是关于第二对角线的对称，故

\[
\dim \mathfrak {g} = 2 l + l ^ {2} + 2 \frac {l (l - 1)}{2} = l (2 l + 1).
\]

令 \(\mathfrak{h}\) 为 \(\mathfrak{g}\) 的对角元素的集合。这是 \(\mathfrak{g}\) 的交换子代数，其基由以下元素组成：

\footnotetext[6]{方阵 \((a_{ij})_{1\leq i,j\leq n}\) 的第二对角线是满足 \(a_{ij}\) 的 \(i+j=n+1\) 所成的族。}

\[
H _ {i} = E _ {i, i} - E _ {- i, - i} \quad (1 \leq i \leq l).
\]

令 \((\varepsilon_{i})\) 为与 \(h^{*}\) 对偶的 \((H_{i})\) 的基。置

\[
\left\{ \begin{array}{l l l} X _ {\varepsilon_ {i}} & = 2 E _ {i, 0} + E _ {0, - i} & (1 \leq i \leq l) \\ X _ {- \varepsilon_ {i}} & = - 2 E _ {- i, 0} - E _ {0, i} & (1 \leq i \leq l) \\ X _ {\varepsilon_ {i} - \varepsilon_ {j}} & = E _ {i, j} - E _ {- j, - i} & (1 \leq i <   j \leq l) \\ X _ {\varepsilon_ {j} - \varepsilon_ {i}} & = - E _ {j, i} + E _ {- i, - j} & (1 \leq i <   j \leq l) \\ X _ {\varepsilon_ {i} + \varepsilon_ {j}} & = E _ {i, - j} - E _ {j, - i} & (1 \leq i <   j \leq l) \\ X _ {- \varepsilon_ {i} - \varepsilon_ {j}} & = - E _ {- j, i} + E _ {- i, j} & (1 \leq i <   j \leq l). \end{array} \right.\tag{3}
\]

容易验证，这些元素构成 \(\mathfrak{h}\) 中 \(\mathfrak{g}\) 的一个补空间的基，并且对于 \(h\in \mathfrak{h}\)，有

\[
[ h, X _ {\alpha} ] = \alpha (h) X _ {\alpha}\tag{4}
\]

对所有 \(\alpha \in \mathbb{R}\) 都成立，其中 \(\mathbb{R}\) 是由 \(\pm \varepsilon_i\) 以及 \(\pm \varepsilon_i \pm \varepsilon_j\)（\(1 \leq i < j \leq l\)）组成的集合。因此，\(\mathfrak{h}\) 等于其在 \(\mathfrak{g}\) 中的正规化子，因而是 \(\mathfrak{g}\) 的嘉当子代数；\(\mathfrak{h}\) 是分裂的，并且 \((\mathfrak{g}, \mathfrak{h})\) 的根正是 \(\mathbb{R}\) 的元素。\(\mathbb{R}\) 的根系 \((\mathfrak{g}, \mathfrak{h})\) 在 \(\mathrm{B}_l\) 时为 \(l \geq 2\) 型，在 \(\mathrm{A}_1\) 时为 \(\mathrm{B}_1\) 型（也称 \(l = 1\) 型）（第 VI 章，§4，第5.I节，推广到 \(l = 1\) 的情形）。因此，\(\mathfrak{g}\) 是 \(\mathrm{B}_l\) 型的可分裂单李代数。

正交代数 \(\mathfrak{o}(\varPsi)\) 的每个分裂嘉当子代数都是 h 在 \(\mathfrak{o}(\varPsi)\) 的某个初等自同构下的变换，因而也是在 \(\mathbf{O}(\varPsi)\) 的某个元素下的变换（参见~(VII)）；所以它就是集合 \(h_{\beta}\)，其中的元素 g 相对于 V 的某个 Witt 基 \(\beta\) 的矩阵为对角矩阵。立即验证可知，在 \(h_{\beta}\) 下不变的唯一向量子空间是由 \(\beta\) 的某个子集生成的子空间。

若 \(l = 1\)，代数 \(\mathfrak{o}(\varPsi)\) 和 \(\mathfrak{sl}(2, k)\) 具有相同的根系，因而同构（另参见~§1，习题16）。下文设 \(l \geq 2\)。

\begin{enumerate}
\def\labelenumi{(\Roman{enumi})}
\setcounter{enumi}{1}
\tightlist
\item
  根系 \(R^{\vee}\) 由第 VI 章，§4，第5.V节确定，我们得到
\end{enumerate}

\[
H _ {\varepsilon_ {i}} = 2 H _ {i}, \quad H _ {\varepsilon_ {i} - \varepsilon_ {j}} = H _ {i} - H _ {j}, \quad H _ {\varepsilon_ {i} + \varepsilon_ {j}} = H _ {i} + H _ {j}.
\]

\begin{enumerate}
\def\labelenumi{(\Roman{enumi})}
\setcounter{enumi}{2}
\tightlist
\item
  置 \(\alpha_{1} = \varepsilon_{1} - \varepsilon_{2},\ldots ,\alpha_{l - 1} = \varepsilon_{l - 1} - \varepsilon_{l},\alpha_{l} = \varepsilon_{l}\)。由第 VI 章，§4，第5.II节，\((\alpha_{1},\dots ,\alpha_{l})\) 是 R 的一个基 B；相对于 B 的正根是 \(\varepsilon_{i}\) 以及 \(\varepsilon_{i}\pm \varepsilon_{j}\)（\(i < j\)）。相应的 Borel 子代数 \(\mathfrak{b}\) 是 \(\mathfrak{g}\) 中上三角矩阵的集合。
\end{enumerate}

立即验证可知，在 b 下稳定且不同于 \(\{0\}\) 和 V 的唯一向量子空间，是与基 \((e_{i})\) 对应的极大旗中的元素，即全迷向子空间 \(V_{1},\ldots,V_{l}\)，其中 \(V_{i}\) 由 \(e_{1},\ldots,e_{i}\) 生成，以及它们的正交补 \(V_{-1},\ldots,V_{-i}\)：\(V_{-i}\) 的正交补 \(V_{i}\) 由 \(e_{1},\ldots,e_{l},e_{0},e_{-l},\ldots,e_{-i-1}\) 生成，且不是全迷向的。另一方面，如果 g 的一个元素使某个向量子空间稳定，它也使其正交补稳定。因此，\(\mathfrak{b}\) 是使旗 \(\mathfrak{g}\) 的各元素稳定的 \(\{\mathrm{V}_1,\dots ,\mathrm{V}_l\}\) 中元素的集合。

如果旗的每个元素都是全迷向的，就称该旗为迷向旗。旗 \(\{V_{1},\ldots,V_{l}\}\) 是极大迷向旗。由于群 \(\mathbf{O}(\varPsi)\) 既在 g 的 Borel 子代数上（参见~(VII)）传递作用，也在极大迷向旗上（《代数》第 IX 章，§4，第3节，定理1）传递作用，我们看到，对于 V 中任意极大迷向旗 \(\delta\)，使 \(b_{\delta}\) 的各元素稳定的 g 中元素的集合 \(\delta\) 是 g 的一个 Borel 子代数，并且映射 \(\delta\mapsto b_{\delta}\) 是从极大迷向旗集合到 g 的 Borel 子代数集合的双射。

令 \(\delta\) 为迷向旗，令 \(\mathfrak{p}_{\delta}\) 为使 \(\mathfrak{g}\) 的各元素稳定的 \(\delta\) 中元素的集合。若 \(\delta \subset \{\mathrm{V}_1, \ldots, \mathrm{V}_l\}\)，则 \(\mathfrak{p}_{\delta}\) 是包含 \(\mathfrak{g}\) 的 \(\mathfrak{b}\) 的一个抛物子代数，并且容易验证，在 \(\neq \{0\}\) 下稳定且不等于 \(\mathfrak{p}_{\delta}\) 的唯一全迷向子空间是 \(\delta\) 的元素。这给出包含 \(2^l\) 的 \(\mathfrak{g}\) 的 \(\mathfrak{b}\) 个抛物子代数。如上所述，映射 \(\delta \mapsto \mathfrak{p}_{\delta}\) 是从 V 中迷向旗的集合到 \(\mathfrak{g}\) 的抛物子代数集合的双射。此外，当且仅当 \(\mathfrak{p}_{\delta} \subset \mathfrak{p}_{\delta'}\) 时，\(\delta \supset \delta'\)。

\begin{enumerate}
\def\labelenumi{(\Roman{enumi})}
\setcounter{enumi}{3}
\tightlist
\item
  与 \(\alpha_{1},\ldots,\alpha_{l}\) 对应的基本权由第 VI 章，§4，第5.VI节给出：
\end{enumerate}

\[
\begin{array}{l} \varpi_ {i} = \varepsilon_ {1} + \dots + \varepsilon_ {i} (1 \leq i \leq l - 1) \\ \varpi_ {l} = \frac {1}{2} (\varepsilon_ {1} + \dots + \varepsilon_ {l}). \end{array}
\]

令 \(\sigma\) 为 \(\mathfrak{g}\) 在 V 上的恒等表示。外幂 \(\bigwedge^{r}\sigma\) 作用于 \(E = \bigwedge^{r}V\)。若 \(h\in \mathfrak{h}\)，

\[
\begin{array}{l} \sigma (h). e _ {i} = \varepsilon_ {i} (h) e _ {i} \quad \text { for } 1 \leq i \leq l \\ \sigma (h). e _ {0} = 0 \\ \sigma (h). e _ {- i} = - \varepsilon_ {i} (h) e _ {- i} \quad \text { for } 1 \leq i \leq l. \end{array}
\]

由此，对于 \(1 \leq r \leq l\)，\(\varepsilon_{1} + \cdots + \varepsilon_{r}\) 是 \(\bigwedge^{r}\sigma\) 的最高权，权为 \(\varepsilon_{1} + \cdots + \varepsilon_{r}\) 的元素正是与 \(e_{1} \wedge \cdots \wedge e_{r}\) 成比例的元素。我们将证明，当 \(1 \leq r \leq l - 1\) 时，表示 \(\bigwedge^{r}\sigma\) 是最高权为 \(\varpi_{r}\) 的基本表示。为此，只需证明 \(\bigwedge^{r}\sigma\) 在 \(0 \leq r \leq 2l + 1\) 时不可约。但是，在 \(\Phi\) 上定义的双线性型 \(\bigwedge^{r}V \times \bigwedge^{2l+1-r}V\) 为

\[
x \wedge y = \varPhi (x, y) e _ {1} \wedge \dots \wedge e _ {l} \wedge e _ {0} \wedge e _ {- l} \wedge \dots \wedge e _ {- 1}
\]

它在 g 下不变，并使 \(\bigwedge^{r}V\) 与 \(\bigwedge^{2l+1-r}V\) 处于对偶关系。因此，表示 \(\bigwedge^{2l+1-r}\sigma\) 是 \(\bigwedge^{r}\sigma\) 的对偶，只需证明 \(\bigwedge^{r}\sigma\) 在 \(0\leq r\leq l\) 时不可约；等价地，只需证明 \(T_{r}\) 中包含 \(\bigwedge^{r}V\) 且在 g 下稳定的最小子空间 \(e_{1}\wedge\cdots\wedge e_{r}\) 就是整个 \(\bigwedge^{r}V\)。当 r=0 和 r=1 时这立即成立（参见~公式（2））。当 r=2（从而 \(l\geq2\)）时，\(\bigwedge^{2}\sigma\) 与 g 的伴随表示（后者不可约）具有相同的维数 \(l(2l + 1)\) 和相同的最高权 \(\varepsilon_1 + \varepsilon_2\)（第 VI 章，§4，第5.IV节）。因此，\(\bigwedge^2\sigma\) 等价于伴随表示，因而不可约。这就证明了 \(l = 1\) 和 \(l = 2\) 时的断言。

现在对 \(l\) 作归纳，并设 \(l \geq r \geq 3\)。先注意，若 W 是 V 的一个奇数维非迷向子空间，正交补为 \(\mathrm{W}'\)，则 \(\Psi_{\mathrm{W}}\)（即 \(\Psi\) 在 W 上的限制）非退化，并且 \(\mathfrak{o}(\Psi_{\mathrm{W}})\) 可与 \(\mathfrak{g}\) 中由在 \(\mathrm{W}'\) 上为零的元素组成的子代数识别。若 \(\dim \mathrm{W} < \dim \mathrm{V}\)，且 \(\Psi_{\mathrm{W}}\) 的指标为最大指标，则归纳假设表明：如果 \(\mathrm{T}_r\) 含有形如 \(w' \wedge w\) 的非零元素，其中 \(w' \in \bigwedge^{r - k} \mathrm{W}'\)，\(w \in \bigwedge^k \mathrm{W}\)（\(0 \leq k \leq r\)），那么 \(\mathrm{T}_r\) 含有 \(w' \wedge \bigwedge^k \mathrm{W}\)：事实上，有 \(a.(w' \wedge w) = w' \wedge a.w\)，对所有 \(a \in \mathfrak{o}(\Psi_{\mathrm{W}})\)。我们对 \(p \in [0, r]\) 作归纳，证明 \(\mathrm{T}_r\) 含有元素

\[
x = e _ {i _ {1}} \wedge \dots \wedge e _ {i _ {r - p}} \wedge e _ {j _ {1}} \wedge \dots \wedge e _ {j _ {p}}
\]

当 \(1 \leq i_{1} < \cdots < i_{r-p} \leq l\) 且 \(-l \leq j_{1} < \cdots < j_{p} \leq 0\) 时。对于 p = 0，这源于 \(\mathfrak{gl}(F)\) 在 \(\bigwedge^{r} F\) 上作用的不可约性（第1节），因为 g 包含使 \(F = V_{l} = \sum_{i=1}^{l} ke_{i}\) 不变并在其上诱导任意自同态的元素（参见~公式（2））。若 p = 1，取 \(q \in (1, l)\)，使得 \(q \neq -j_{1}\)，且存在 \(\lambda \in [1, r - p]\) 使 \(q = i_{\lambda}\)；若 \(p \geq 2\)，取 \(q \in [1, l]\)，使得 \(-q \in \{j_{1}, \ldots, j_{p}\}\)。必要时置换 \(e_{i}\)，可设 q = 1。现在取 W 为 \(W' = ke_{1} + ke_{-1}\) 的正交补。若 p = 1，则有 \(x \in e_{1} \wedge \bigwedge^{r-1} W\)；由于 \(T_{r}\) 含有 \(e_{1} \wedge \cdots \wedge e_{r}\)，我们看到 \(T_{r}\) 含有 x。若 \(p \geq 2\)，则或者 \(x \in e_{-1} \wedge \bigwedge^{r-1} W\)，或者 \(x \in e_{1} \wedge e_{-1} \wedge \bigwedge^{r-2} W\)；由于根据归纳假设，\(T_{r}\) 含有 \(e_{-1} \wedge e_{2} \wedge \cdots \wedge e_{r-1}\) 和 \(e_{-1} \wedge e_{1} \wedge e_{2} \wedge \cdots \wedge e_{r-2}\)，我们看到 \(T_{r}\) 含有 x，从而证明完毕。

关于 \(\bigwedge^{r}\sigma\) 不可约性的另一证明，见习题6。

现在确定最高权为 \(\varpi_{l}\) 的基本表示。

引理 1. 设 \(\mathrm{V}\) 为有限维向量空间，\(\mathbb{Q}\) 为 \(\mathrm{V}\) 上的非退化二次型，\(\Psi\) 为与 \(\mathbb{Q}\) 相关的对称双线性型，\(\mathrm{C(Q)}\) 为相对于 \(\mathrm{V}\) 的 \(\mathbb{Q}\) 的 Clifford 代数，\(f_0\) 为以下典则映射的复合

\[
\mathfrak {o} (\Psi) \longrightarrow \mathfrak {g l} (V) \longrightarrow V \otimes V ^ {*} \longrightarrow V \otimes V \longrightarrow C ^ {+} (Q)
\]

（第一个是典则嵌入，第三个由对应于 \(\mathrm{V}^*\) 的从 \(\mathrm{V}\) 到 \(\Psi\) 的典则同构定义，第四个由 \(\mathrm{C(Q)}\) 中的乘法定义，参见~《代数》第 IX 章，§ 9，第1节）。置 \(f = \frac{1}{2} f_0\)。

\begin{enumerate}
\def\labelenumi{(\roman{enumi})}
\item
  若 \((e_r), (e_r')\) 是 V 的满足 \(\Psi(e_r, e_s') = \delta_{rs}\) 的基，则对所有 \(f_0(a) = \sum_r (ae_r)e_r'\)，有 \(a \in \mathfrak{o}(\Psi)\)。
\item
  若 \(a, b \in \mathfrak{o}(\Psi)\)，则有 \(\sum_{r}(ae_{r})(be_{r}^{\prime}) = -\sum_{r}(abe_{r})e_{r}^{\prime}\)。
\item
  若 \(a \in \mathfrak{o}(\varPsi)\) 且 \(v \in V\)，则有 \([f(a), v] = av\)。
\item
  若 \(a, b \in \mathfrak{o}(\varPsi)\)，则有 \([f(a), f(b)] = f([a, b])\)。
\item
  \(f(\mathfrak{o}(\Psi))\) 生成结合代数 \(\mathrm{C}^{+}(\mathrm{Q})\)。
\item
  设 \(\mathrm{N}\) 为左 \(\mathrm{C}^{+}(\mathrm{Q})\)-模，\(\rho\) 为从 \(\mathrm{C}^{+}(\mathrm{Q})\) 到 \(\operatorname{End}_k(\mathrm{N})\) 的相应同态。则 \(\rho \circ f\) 是 \(\mathfrak{o}(\Psi)\) 在 \(\mathrm{N}\) 上的表示。若 \(\mathrm{N}\) 是单模，则 \(\rho \circ f\) 不可约。
\end{enumerate}

断言 (i) 是显然的。若 \(a, b \in \mathfrak{o}(\varPsi)\)，记 \(\varPsi(x,y) = \langle x,y\rangle\)，则有：

\[
\begin{array}{c} \sum_ {r} (a e _ {r}) (b e _ {r} ^ {\prime}) = \sum_ {r, s, t} \langle a e _ {r}, e _ {s} ^ {\prime} \rangle \langle b e _ {r} ^ {\prime}, e _ {t} \rangle e _ {s} e _ {t} ^ {\prime} = \sum_ {r, s, t} \langle e _ {r}, a e _ {s} ^ {\prime} \rangle \langle e _ {r} ^ {\prime}, b e _ {t} \rangle e _ {s} e _ {t} ^ {\prime} \\ = \sum_ {s, t} \langle a e _ {s} ^ {\prime}, b e _ {t} \rangle e _ {s} e _ {t} ^ {\prime} = - \sum_ {s, t} \langle e _ {s} ^ {\prime}, a b e _ {t} \rangle e _ {s} e _ {t} ^ {\prime} = - \sum_ {t} (a b e _ {t}) e _ {t} ^ {\prime} \end{array}
\]

这证明了 (ii)。接着，对所有 \(v \in V\)，由 (i) 得

\[
\begin{array}{l} [ f (a), v ] = \frac {1}{2} \sum_ {r} ((a e _ {r}) e _ {r} ^ {\prime} v - v (a e _ {r}) e _ {r} ^ {\prime}) \\ \qquad = \frac {1}{2} \sum_ {r} ((a e _ {r}) e _ {r} ^ {\prime} v + (a e _ {r}) v e _ {r} ^ {\prime} - (a e _ {r}) v e _ {r} ^ {\prime} - v (a e _ {r}) e _ {r} ^ {\prime}) \\ \qquad = \frac {1}{2} \sum_ {r} ((a e _ {r}) \langle e _ {r} ^ {\prime}, v \rangle - \langle a e _ {r}, v \rangle e _ {r} ^ {\prime}) \\ \qquad = \frac {1}{2} a \left(\sum_ {r} \langle e _ {r} ^ {\prime}, v \rangle e _ {r}\right) + \frac {1}{2} \sum_ {r} \langle e _ {r}, a v \rangle e _ {r} ^ {\prime} = \frac {1}{2} a v + \frac {1}{2} a v = a v, \end{array}
\]

这证明了 (iii)。于是

\[
\begin{array}{l} [ f (a), f (b) ] = \left[ f (a), \frac {1}{2} \sum_ {r} (b e _ {r}) e _ {r} ^ {\prime} \right] \quad \text { by   (i) } \\ \quad = \frac {1}{2} \sum_ {r} ([ f (a), b e _ {r} ] e _ {r} ^ {\prime} + (b e _ {r}) [ f (a), e _ {r} ^ {\prime} ]) \\ \quad = \frac {1}{2} \sum_ {r} ((a b e _ {r}) e _ {r} ^ {\prime} + (b e _ {r}) (a e _ {r} ^ {\prime})) \quad \text { by   (iii) } \\ \quad = \frac {1}{2} \sum_ {r} ((a b e _ {r}) e _ {r} ^ {\prime} - (b a e _ {r}) e _ {r} ^ {\prime}) \quad \text { by   (ii) } \\ \quad = f ([ a, b ]) \quad \text { by   (i) } \end{array}
\]

这证明了 (iv)。为证明 (v)，通过扩张标量，可设 \(k\) 代数闭。于是取 V 的一个基 \((e_r)\)，使 \(\varPsi(e_r, e_s) = \delta_{rs}\)，从而 \(e_r' = e_r\)。若 \(i \neq j\)，则 \(\mathrm{E}_{ij} - \mathrm{E}_{ji} \in \mathfrak{o}(\varPsi)\)，且

\[
f (\mathrm{E} _ {i j} - \mathrm{E} _ {j i}) = \frac {1}{2} (e _ {i} e _ {j} - e _ {j} e _ {i}) = e _ {i} e _ {j};
\]

但 \(e_{i}e_{j}\) 生成 \(\mathrm{C}^{+}(\mathrm{Q})\)。

断言 (vi) 由 (iv) 和 (v) 得出。

Q.E.D.

现在回顾相关记号。令 \(\tilde{V} = F + F'\)、\(\tilde{Q}\)（分别为 \(\tilde{\Psi}\)）表示 Q（分别为 \(\Psi\)）在 \(\tilde{V}\) 上的限制。于是，\(\tilde{Q}\) 是定义在 \(\tilde{V}\) 上、维数为 2l 且指标为最大指标 l 的非退化二次型，而 Clifford 代数 \(C(\tilde{Q})\) 是维数为 \(2^{2l}\) 的中心单代数。令 N 为极大迷向子空间 \(F'\) 的外代数，它由 \(e_{-1}, \ldots, e_{-l}\) 生成。利用 \(F'\) 将 F 识别为其对偶空间，并借助 \(\Psi\)，对于 \(x \in F'\)（分别对于 \(y \in F\)），在 N 中将以 x（分别以 y）作左外积（分别作左内积）记为 \(\lambda(x)\)（分别记为 \(\lambda(y)\)）；若 \(a_{1}, \ldots, a_{k} \in F'\)，则

\[
\lambda (x). (a _ {1} \wedge \dots \wedge a _ {k}) = x \wedge a _ {1} \wedge \dots \wedge a _ {k}
\]

\[
\lambda (y). \left(a _ {1} \wedge \dots \wedge a _ {k}\right) = \sum_ {i = 1} ^ {k} (- 1) ^ {i - 1} \Psi \left(a _ {i}, y\right) a _ {1} \wedge \dots \wedge a _ {i - 1} \wedge a _ {i + 1} \wedge \dots \wedge a _ {k}.
\]

容易验证 \(\lambda (x)^{2} = \lambda (y)^{2} = 0\)，并且

\[
\lambda (x) \lambda (y) + \lambda (y) \lambda (x) = \Psi (x, y). 1.
\]

由此（《代数》第 IX 章，§9，第1节），存在唯一的同态（仍记为 \(\lambda\)）从 \(\mathrm{C}(\tilde{\mathbf{Q}})\) 到 \(\mathrm{End}(\mathrm{N})\)，延拓映射 \(\lambda : \mathrm{F} \cup \mathrm{F}' \to \mathrm{End}(\mathrm{N})\)。由于 \(\dim \mathrm{N} = 2^l\)，且 \(\mathrm{C}(\tilde{\mathbf{Q}})\) 只有一类维数为 \(2^l\) 的单模（《代数》第 IX 章，§9，第4节，定理2），由 \(\mathrm{C}(\tilde{\mathbf{Q}})\) 定义的 \(\mathrm{N}\) 在 \(\lambda\) 上的表示是不可约表示，并且是 \(\mathrm{C}(\tilde{\mathbf{Q}})\) 的旋量表示（同上引）。

现在考虑从 \(\mu:v\mapsto e_{0}v\) 到 \(\tilde{\mathrm{V}}\) 的映射 \(\mathrm{C}^{+}(\mathrm{Q})\)。对于 \(v\in\tilde{V}\)，有

\[
(e _ {0} v) ^ {2} = - e _ {0} ^ {2} v ^ {2} = - \mathrm{Q} (e _ {0}) \mathrm{Q} (v) = \mathrm{Q} (v) = \tilde {\mathrm{Q}} (v)
\]

并且 \(\mu\) 唯一延拓为同态（仍记为 \(\mu\)），从 \(\mathrm{C}(\tilde{\mathbb{Q}})\) 到 \(\mathrm{C}^{+}(\mathrm{Q})\)。由于 \(\mathrm{C}(\tilde{\mathbb{Q}})\) 是单代数，且

\[
\dim \mathrm{C} ^ {+} (\mathrm{Q}) = \dim \mathrm{C} (\tilde {\mathrm{Q}}) = 2 ^ {2 l},
\]

可知 \(\mu\) 是同构。因此，\(\lambda \circ \mu^{-1}\) 在 N 上定义了一个单 \(\mathrm{C}^{+}(\mathrm{Q})\)-模结构，而 \(\rho = \lambda \circ \mu^{-1} \circ f\) 是 g 在 N 上的不可约表示（引理1 (vi)）。

另一方面，根据引理1 (i)，有

\[
f (H _ {i}) = \frac {1}{2} (e _ {i} e _ {- i} - e _ {- i} e _ {i}).
\]

由于 \(e_{i}e_{-i} = -e_{0}^{2}e_{i}e_{-i} = e_{0}e_{i}e_{0}e_{-i}\)，且 \(e_{i}e_{-i} + e_{-i}e_{i} = 1\)，有

\[
\mu^ {- 1} \circ f (H _ {i}) = \frac {1}{2} - e _ {- i} e _ {i} = - \frac {1}{2} + e _ {i} e _ {- i}.
\]

由此，对于 \(1 \leq i_{1} < \cdots < i_{k} \leq l\)，有：

\[
\rho (H _ {i}) (e _ {- i _ {1}} \wedge \dots \wedge e _ {- i _ {k}}) = \left\{ \begin{array}{l l} - \frac {1}{2} e _ {- i _ {1}} \wedge \dots \wedge e _ {- i _ {k}} & \text { if } i \in \{i _ {1}, \ldots , i _ {k} \} \\ \frac {1}{2} e _ {- i _ {1}} \wedge \dots \wedge e _ {- i _ {k}} & \text { if } i \notin \{i _ {1}, \ldots , i _ {k} \} \end{array} \right.
\]

并且对于 \(h\in \mathfrak{h}\)

\[
\begin{array}{l} \rho (h) (e _ {- i _ {1}} \wedge \dots \wedge e _ {- i _ {k}}) \\ = (\frac {1}{2} (\varepsilon_ {1} + \dots + \varepsilon_ {l}) - (\varepsilon_ {i _ {1}} + \dots + \varepsilon_ {i _ {k}})) (h) (e _ {- i _ {1}} \wedge \dots \wedge e _ {- i _ {k}}). \end{array} \tag {5}
\]

这表明 \(\rho\) 的最高权是 \(\varpi_{l}\)。我们称 \(\rho\) 为 g 的旋量表示。注意，其权全为单权（此外，\(\varpi_{l}\) 是极小权）。

\begin{enumerate}
\def\labelenumi{(\Alph{enumi})}
\setcounter{enumi}{21}
\tightlist
\item
有 \(w_0 = -1\)，所以 \(\mathfrak{g}\) 的每个有限维单表示都是正交型或辛型。由第 VI 章的结果，\(\varpi_r\) 关于 \((\alpha_1, \ldots, \alpha_l)\) 的坐标之和在 \(1 \leq r \leq l - 1\) 时为整数；因此，表示 \(\bigwedge^r \sigma\) 是正交型的。此外，它保持 \(\Psi_{(r)}\)（\(\Psi\) 在 \(\bigwedge^r V\) 上的延拓）不变。
\end{enumerate}

对于旋量表示，\(\varpi_{l}\) 关于 \((\alpha_{1},\ldots,\alpha_{l})\) 的坐标之和为 \(\frac{1}{2}(1+\cdots+l)=\frac{l(l+1)}{4}\)（同上引）。因此，当 \(l\equiv0\) 或 -1 (mod. 4) 时它是正交型的，当 \(l\equiv1\) 或 2 (mod. 4) 时它是辛型的。事实上，考虑 \(\Phi\) 上的双线性型 \(N=\bigwedge F'\)，定义如下：若 \(x\in\bigwedge^{p}F'\) 且 \(y\in\bigwedge^{q}F'\)，则当 \(\Phi(x,y)=0\) 时置 \(p+q\neq l\)，并且

\[
x \wedge y = (- 1) ^ {\frac {p (p + 1)}{2}} \varPhi (x, y) e _ {- 1} \wedge \dots \wedge e _ {- r}
\]

当 \(p + q = l\) 时，容易验证 \(\Phi\) 非退化；当 \(l \equiv 0, -1 \pmod{4}\) 时，\(l \equiv 1, 2 \pmod{4}\) 也是非退化的。

\[
f (X _ {\varepsilon_ {i}}) = e _ {0} e _ {i}, \quad f (X _ {- \varepsilon_ {i}}) = - e _ {0} e _ {- i}
\]

对 \(1 \leq i \leq l\)，且

\[
f (X _ {\varepsilon_ {i} - \varepsilon_ {j}}) = \frac {1}{2} (e _ {i} e _ {- j} - e _ {- j} e _ {i}) = e _ {i} e _ {- j} = e _ {0} e _ {i} e _ {0} e _ {- j}
\]

对于 \(1 \leq i < j \leq l\)，同理

\[
f (X _ {\varepsilon_ {j} - \varepsilon_ {i}}) = - e _ {0} e _ {j} e _ {0} e _ {- i}, f (X _ {\varepsilon_ {i} + \varepsilon_ {j}}) = e _ {0} e _ {i} e _ {0} e _ {j}, f (X _ {- \varepsilon_ {i} - \varepsilon_ {j}}) = e _ {0} e _ {- i} e _ {0} e _ {- j};
\]

从而

\[
\mu^ {- 1} \circ f (X _ {\varepsilon_ {i}}) = e _ {i}, \quad \mu^ {- 1} f (X _ {- \varepsilon_ {i}}) = - e _ {- i}
\]

且

\[
\mu^ {- 1} \circ f (X _ {\pm \varepsilon_ {i} \pm \varepsilon_ {j}}) = c e _ {\pm i} e _ {\pm j} \quad \text {   for   } 1 \leq i, j \leq l, i \neq j \text {   with   } c \in \{1, - 1 \}.
\]

现在很容易验证，\(\Phi\) 实际上在 g 下不变（参见~习题18）。

\begin{enumerate}
\def\labelenumi{(\Roman{enumi})}
\setcounter{enumi}{5}
\tightlist
\item
对于 \(x \in \mathfrak{g}\)，\(\sigma(x)\) 的特征多项式具有形式
\end{enumerate}

\[
\mathrm{T} ^ {2 l + 1} + f _ {1} (x) \mathrm{T} ^ {2 l} + f _ {2} (x) \mathrm{T} ^ {2 l - 1} + \dots + f _ {2 l + 1} (x)
\]

其中 \(f_{1},\ldots ,f_{2l + 1}\) 是 \(\mathfrak{g}\) 上的不变多项式函数。

若 \(x = \xi_{1}H_{1} + \cdots + \xi_{l}H_{l} \in h\)，则 \(f_{i}(x)\) 在相差一个符号的意义下是 \(\xi_{1}, \ldots, \xi_{l}, -\xi_{1}, \ldots, -\xi_{l}\) 的初等对称函数；这些对称函数在奇数次数时为零，并且

\[
\mathrm{T} ^ {2 l + 1} + f _ {2} (x) \mathrm{T} ^ {2 l - 1} + f _ {4} (x) \mathrm{T} ^ {2 l - 3} + \dots + f _ {2 l} (x) \mathrm{T} = \mathrm{T} (\mathrm{T} ^ {2} - \xi_ {1} ^ {2}) \dots (\mathrm{T} ^ {2} - \xi_ {l} ^ {2})
\]

从而 \(f_{2},\ldots,f_{2l}\) 在相差一个符号的意义下是 \(\xi_{1}^{2},\ldots,\xi_{l}^{2}\) 的初等对称函数，它们是 \(\mathbf{S}(\mathfrak{h}^{*})^{\mathrm{W}}\) 的代数无关生成元（第 VI 章，§4，第5.IX节）。根据 §8，第3节，定理1 (i)，可知 \(f_{1}=f_{3}=f_{5}=\cdots=0\)，并且 \((f_{2},f_{4},\ldots,f_{2l})\) 是生成 g 上不变多项式函数代数的代数自由族。

\begin{enumerate}
\def\labelenumi{(\Roman{enumi})}
\setcounter{enumi}{6}
\tightlist
\item
由于邓金图的唯一自同构是恒等自同构，有 \(\mathrm{Aut}(\mathfrak{g}) = \mathrm{Aut}_0(\mathfrak{g})\)。
\end{enumerate}

令 \(\Sigma\) 为关于 \(\Psi\) 的 V 的相似变换群。对所有 \(g \in \Sigma\)，令 \(\varphi(g)\) 为自同构 \(x \mapsto gxg^{-1}\)，它作用于 \(\mathfrak{g}\)。于是 \(\varphi\) 是从 \(\Sigma\) 到 \(\operatorname{Aut}(\mathfrak{g})\) 的同态。我们证明它是满射。令 \(\alpha \in \operatorname{Aut}(\mathfrak{g}) = \operatorname{Aut}_0(\mathfrak{g})\)。由 §7，第1节，命题2，存在 \(s \in \mathbf{GL}(V)\)，使得 \(\alpha(x) = sxs^{-1}\) 对所有 \(x \in \mathfrak{g}\) 成立。于是 \(s\) 将 \(\Psi\) 变为双线性型 \(\Psi'\)，它在 \(\mathfrak{g}\) 下不变，从而与 \(\Psi\) 成比例；这证明 \(s \in \Sigma\)。

由于 g 的恒等表示不可约，其交换子仅为标量（\(\S6\)，第1节，命题1），所以 \(\varphi\) 的核为 \(k^{*}\)。因此，群 \(\operatorname{Aut}(\mathfrak{g}) = \operatorname{Aut}_{0}(\mathfrak{g})\) 可与 \(\Sigma/k^{*}\) 识别。但是，由《代数》第 IX 章，\(\S6\)，第5节，\(\Sigma\) 是 \(k^{*}\) 与 \(\mathbf{SO}(\Psi)\) 的乘积；所以 \(\operatorname{Aut}(\mathfrak{g}) = \operatorname{Aut}_{0}(\mathfrak{g})\) 可与 \(\mathbf{SO}(\Psi)\) 识别。

令 \(\mathbf{O}_0^+(\varPsi)\) 为 \(\varPsi\) 的既约正交群（《代数》第 IX 章，§9，第5节）。由于 \(\mathbf{SO}(\varPsi)/\mathbf{O}_0^+(\varPsi)\) 是交换群（同上引），群 \(\mathrm{Aut}_e(\mathfrak{g})\) 包含于 \(\mathbf{O}_0^+(\varPsi)\)（§11，第2节，命题3）；事实上它等于该群（习题7）。

\begin{enumerate}
\def\labelenumi{(\Roman{enumi})}
\setcounter{enumi}{7}
\tightlist
\item
\(\Phi_{R}\) 上的典则双线性型 \(h^{*}\) 为
\end{enumerate}

\[
\varPhi_ {\mathrm{R}} (\xi_ {1} \varepsilon_ {1} + \dots + \xi_ {l} \varepsilon_ {l}, \xi_ {1} ^ {\prime} \varepsilon_ {1} + \dots + \xi_ {l} ^ {\prime} \varepsilon_ {l}) = \frac {1}{4 l - 2} (\xi_ {1} \xi_ {1} ^ {\prime} + \dots + \xi_ {l} \xi_ {l} ^ {\prime})
\]

（第 VI 章，§4，第5.V节）。由 \(\mathfrak{h}\) 定义的从 \(\mathfrak{h}^*\) 到 \(\varPhi_{\mathrm{R}}\) 的同构将 \(H_{i}\) 映为 \((4l - 2).\varepsilon_{i}\)。因此，\(\varPhi_{\mathrm{R}}\) 的逆型，即 Killing 型在 \(\mathfrak{h}\) 上的限制，为

\[
\varPhi (\xi_ {1} H _ {1} + \dots + \xi_ {l} H _ {l}, \xi_ {1} ^ {\prime} H _ {1} + \dots + \xi_ {l} ^ {\prime} H _ {l}) = (4 l - 2) (\xi_ {1} \xi_ {1} ^ {\prime} + \dots + \xi_ {l} \xi_ {l} ^ {\prime}).
\]

\begin{enumerate}
\def\labelenumi{(\Roman{enumi})}
\setcounter{enumi}{8}
\tightlist
\item
回顾由公式（3）定义的 \(X_{\alpha}\)（\(\alpha \in R\)）。容易验证，对于 \([X_{\alpha}, X_{-\alpha}] = -H_{\alpha}\)，有 \(\alpha \in R\)。另一方面，令 M 为矩阵 \(I + E_{0,0}\)；由于 \(M = S^{t}M^{-1}S\)，映射
\end{enumerate}

\[
\theta : g \mapsto - M ^ {- 1 t} g M
\]

是 \(\mathfrak{g}\) 的自同构，且 \(\theta(X_{\alpha}) = X_{-\alpha}\) 对所有 \(\alpha \in \mathrm{R}\) 成立。因此，\((X_{\alpha})\) 是 \((\mathfrak{g},\mathfrak{h})\) 中的谢瓦莱系。

假设 \(k = \mathbf{Q}\)。嘉当子代数 \(\mathfrak{h}\) 有两个许可格：格 \(\mathrm{Q}(\mathrm{R}^{\vee})\) 由 \(H_{\alpha}\) 生成，格 \(\mathrm{P}(\mathrm{R}^{\vee})\) 由 \(H_{i}\) 生成，且由 \(\mathfrak{h}\) 中的整数项对角矩阵组成。因此，\(\mathfrak{o}_S(2l + 1, \mathbf{Z})\)（即 \(\mathfrak{g}\) 中具有整数项的矩阵集合）是谢瓦莱序 \(\mathrm{P}(\mathrm{R}^{\vee}) + \sum \mathbf{Z}.X_{\alpha}\)（属于 \(\mathfrak{g}\)）。由于 \((X_{\pm \varepsilon_i})^2 = 2E_{\pm i,\mp i}, (X_{\pm \varepsilon_i})^3 = 0\)，且 \((X_{\pm \varepsilon_i\pm \varepsilon_j})^2 = 0\)，可知格 \(\mathcal{V}\) 由 Witt 基 \((e_i)_{-l\leq i\leq l}\) 生成，是 \(\mathfrak{o}_S(2l + 1,\mathbf{Z})\) 中关于 \(V\) 的容许格。\(\bigwedge^r\mathcal{V}\) 在 \(\bigwedge^r V\) 中也如此。

现在考虑 g 在 \(\rho\) 上的旋量表示 \(N = \bigwedge F'\)。由于其各权不将 \(\mathrm{P}(\mathrm{R}^{\vee})\) 映入 Z，它对 \(\mathfrak{o}_{S}(2l + 1, \mathbf{Z})\) 没有容许格。另一方面，由 N 的典则基 \((e_{-i_{1}} \wedge \cdots \wedge e_{-i_{k}})\)（\(1 \leq i_{1} < \cdots < i_{k} \leq l\)）生成的格 N 是谢瓦莱序 \(\mathcal{G} = \mathrm{Q}(\mathrm{R}^{\vee}) + \sum_{\alpha \in \mathrm{R}} \mathbf{Z}. X_{\alpha}\) 的容许格。事实上，显然 N 在以 \(e_{-i}\) 作外积以及以 \(e_{i}\) 作内积下稳定（\(1 \leq i \leq l\)）。(V) 中的公式于是表明 N 在 \(\rho(\mathcal{G})\) 下稳定。此外，由于对所有 \(\rho(X_{\alpha})^{2} = 0\) 都有 \(\alpha \in R\)，可知 N 是容许格。

\subsection{\texorpdfstring{类型为 \(C_{l}\) 的代数 ( \(l \geq 1\) )}{类型为 C\_\{l\} 的代数 ( l \textbackslash geq 1 )}}

\begin{enumerate}
\def\labelenumi{(\Roman{enumi})}
\tightlist
\item
设 \(\varPsi\) 为定义在有限维向量空间 \(V\)（维数 \(2l \geq 2\)）上的非退化交错双线性型；\(x\) 的自同态 \(V\) 若满足 \(\varPsi(xv, v') + \varPsi(v, xv') = 0\)（对所有 \(v, v' \in V\)），则这些自同态构成 \(\mathfrak{sl}(V)\) 的半单李子代数，记为 \(\mathfrak{sp}(\varPsi)\)，称其为与 \(\varPsi\) 相关的辛李代数。
\end{enumerate}

由《代数》第 IX 章，§4，第2节，V 可写成两个极大全迷向子空间 F 和 \(F'\) 的直和，它们关于 \(\Psi\) 处于对偶关系。令 \((e_{i})_{1\leq i\leq l}\) 为 F 的一个基，令 \((e_{-i})_{1\leq i\leq l}\) 为 \(F'\) 的对偶基。于是

\[
(e _ {1}, \dots , e _ {l}, e _ {- l}, \dots , e _ {- 1})
\]

是 V 的一个基；我们称其为 V 的 Witt 基（或辛基）。相对于此基，\(\varPsi\) 的矩阵是阶为 2l 的方阵

\[
J = \left( \begin{array}{c c} 0 & s \\ - s & 0 \end{array} \right)
\]

其中 s 是阶为 l 的方阵，除第二对角线上的元素等于 1 外，其余元素全为 0，参见~第2.I节。

代数 \(\mathfrak{g} = \mathfrak{sp}(\varPsi)\) 可与方阵代数 \(\mathfrak{sp}(2l,k)\) 识别，后者由方阵 \(a\) 构成，方阵的阶为 \(2l\)，且满足 \(a = -J^{-1t}aJ = J^t aJ\)。

\[
a = \left( \begin{array}{c c} A & B \\ C & - s ^ {t} A s \end{array} \right)
\]

其中 A、B、C 是阶为 l 的方阵，并满足 \(B = s^{t}Bs\) 且 \(C = c^{t}Cs\)；换言之，B 和 C 关于第二对角线对称。因此

\[
\dim \mathfrak {g} = l ^ {2} + 2 \frac {l (l + 1)}{2} = l (2 l + 1).
\]

令 \(\mathfrak{h}\) 为 \(\mathfrak{g}\) 中对角矩阵的集合。这是 \(\mathfrak{g}\) 的交换子代数，其基由元素 \(H_{i} = E_{i,i} - E_{-i,-i}\)（\(1 \leq i \leq l\)）组成。令 \((\varepsilon_{i})_{1 \leq i \leq l}\) 为 \((H_{i})\) 的对偶基。对于 \(1 \leq i < j \leq l\)，置

\[
\left\{ \begin{array}{l l} X _ {2 \varepsilon_ {i}} & = E _ {i, - i} \\ X _ {- 2 \varepsilon_ {i}} & = - E _ {- i, i} \\ X _ {\varepsilon_ {i} - \varepsilon_ {j}} & = E _ {i, j} - E _ {- j, - i} \\ X _ {- \varepsilon_ {i} + \varepsilon_ {j}} & = - E _ {j, i} + E _ {- i, - j} \\ X _ {\varepsilon_ {i} + \varepsilon_ {j}} & = E _ {i, - j} + E _ {j, - i} \\ X _ {- \varepsilon_ {i} - \varepsilon_ {j}} & = - E _ {- i, j} - E _ {- j, i}. \end{array} \right.\tag{6}
\]

容易验证，这些元素构成 \(\mathfrak{h}\) 中 \(\mathfrak{g}\) 的一个补空间的基，并且对于 \(h\in \mathfrak{h}\)，有

\[
[ h, X _ {\alpha} ] = \alpha (h) X _ {\alpha}\tag{7}
\]

对所有 \(\alpha \in \mathbb{R}\) 都成立，其中 \(\mathbb{R}\) 是由 \(\pm 2\varepsilon_{i}\) 以及 \(\pm \varepsilon_{i} \pm \varepsilon_{j}\)（\(i < j\)）组成的集合。因此，\(\mathfrak{h}\) 等于其在 \(\mathfrak{g}\) 中的正规化子，因而是 \(\mathfrak{g}\) 的嘉当子代数；\(\mathfrak{h}\) 是分裂的，并且 \((\mathfrak{g}, \mathfrak{h})\) 的根正是 \(\mathbb{R}\) 的元素。\(\mathbb{R}\) 的根系 \((\mathfrak{g}, \mathfrak{h})\) 在 \(\mathrm{C}_{l}\) 时为 \(l \geq 2\) 型，在 \(\mathrm{A}_{1}\) 时为 \(\mathrm{C}_{1}\) 型（即 \(l = 1\) 型）（第 VI 章，§4，第6.I节，推广到 \(l = 1\) 的情形）。因此，\(\mathfrak{g}\) 是 \(\mathrm{C}_{l}\) 型的可分裂单李代数。

g 的每个分裂嘉当子代数都可由一个初等自同构变换为 h，因而也可由辛群 \(\mathbf{Sp}(\varPsi)\) 的某个元素变换得到（参见~(VII)）；所以它是集合 \(h_{\beta}\)，其中的元素 g 相对于 V 的某个 Witt 基 \(\beta\) 的矩阵为对角矩阵。立即验证可知，在 \(h_{\beta}\) 下不变的唯一向量子空间是由 \(\beta\) 的某个子集生成的子空间。

有 \(\mathfrak{sp}(2,k) = \mathfrak{sl}(2,k)\)。另一方面，代数 \(\mathfrak{sp}(4,k)\) 和 \(\mathfrak{o}_S(5,k)\) 具有相同的根系，因而同构（另参见~习题3）。下文设 \(l \geq 2\)。

\begin{enumerate}
\def\labelenumi{(\Roman{enumi})}
\setcounter{enumi}{1}
\tightlist
\item
根系 \(R^{\vee}\) 由第 VI 章，§4，第6.I节和第6.V节确定；我们得到
\end{enumerate}

\[
H _ {2 \varepsilon_ {i}} = H _ {i}, H _ {\varepsilon_ {i} - \varepsilon_ {j}} = H _ {i} - H _ {j}, H _ {\varepsilon_ {i} + \varepsilon_ {j}} = H _ {i} + H _ {j}.
\]

\begin{enumerate}
\def\labelenumi{(\Roman{enumi})}
\setcounter{enumi}{2}
\tightlist
\item
置 \(\alpha_{1} = \varepsilon_{1} - \varepsilon_{2},\ldots ,\alpha_{l - 1} = \varepsilon_{l - 1} - \varepsilon_{l},\alpha_{l} = 2\varepsilon_{l}\)。由第 VI 章，§4，第6.II节，\(\{\alpha_1,\dots ,\alpha_l\}\) 是 R 的一个基 B；相对于 B 的正根是 \(2\varepsilon_{i}\) 以及 \(\varepsilon_{i}\pm \varepsilon_{j}\)（\(i < j\)）。相应的 Borel 子代数 \(\mathfrak{b}\) 是 \(\mathfrak{g}\) 中上三角矩阵的集合。
\end{enumerate}

令 \(\delta\) 为 V 中的迷向旗（即其元素都是关于 \(\varPsi\) 的全迷向子空间），令 \(\mathfrak{p}_{\delta}\) 为由使 \(\mathfrak{g}\) 的各元素稳定的 \(\delta\) 中元素组成的子代数。像第2.III节中那样，我们证明映射 \(\delta \mapsto \mathfrak{p}_{\delta}\) 是从迷向旗的集合（分别从极大迷向旗的集合）到 \(\mathfrak{g}\) 的抛物子代数集合（分别到 Borel 子代数集合）的双射；当且仅当 \(\mathfrak{p}_{\delta} \supset \mathfrak{p}_{\delta'}\) 时，有 \(\delta \subset \delta'\)。

\begin{enumerate}
\def\labelenumi{(\Roman{enumi})}
\setcounter{enumi}{3}
\tightlist
\item
与 \(\alpha_{1},\ldots,\alpha_{l}\) 对应的基本权由第 VI 章，§4，第6.VI节给出，即 \(\varpi_{i}=\varepsilon_{1}+\cdots+\varepsilon_{i}\ (1\leq i\leq l)\)。
\end{enumerate}

我们将说明，权为 \(\sigma_{r}\) 的基本表示 \(\varpi_{r}\) 可以实现为 \(\bigwedge^{r}\sigma\) 的一个子表示，其中 \(\sigma\) 是 g 在 V 上的恒等表示；为此，我们将研究 g 在外代数 \(\bigwedge\sigma\) 上的表示 \(\bigwedge V\) 的分解。

令 \((e_{i}^{*})\) 为 \(V^{*}\) 中与 \((e_{i})\) 对偶的基。交错双线性型 \(\varPsi\) 可识别为 \(\varGamma^{*}\in\bigwedge^{2}V^{*}\) 中的元素，容易验证

\[
\Gamma^ {*} = - \sum_ {i = 1} ^ {l} e _ {i} ^ {*} \wedge e _ {- i} ^ {*}.
\]

令 \(\varPsi^{*}\) 为 \(\varPsi\) 的逆型（《代数》第 IX 章，§1，第7节）；立即有

\[
\varPsi^ {*} (e _ {i} ^ {*}, e _ {j} ^ {*}) = 0
\]

当 \(i \neq -j\) 时，\(\Psi^{*}(e_{i}^{*}, e_{-i}^{*}) = -1\)；且对 \(1 \leq i \leq l\)，\(\Psi^{*}\)。若将 \(\Gamma \in \bigwedge^{2}\mathrm{V}\) 视为相应的元素，则

\[
\Gamma = \sum_ {i = 1} ^ {l} e _ {i} \wedge e _ {- i}.
\]

记 \(X_{-}\) 为 \(\bigwedge V\) 上由以 \(\Gamma\) 作左外积给出的自同态，记 \(X_{+}\) 为 \(\bigwedge V\) 上由以 \(-\Gamma^{*}\) 作左内积给出的自同态：

\[
X _ {-} u = \left(\sum_ {i = 1} ^ {l} e _ {i} \wedge e _ {- i}\right) \wedge u,
\]

\[
X _ {+} u = \left(\sum_ {i = 1} ^ {l} e _ {i} ^ {*} \wedge e _ {- i} ^ {*}\right) \lrcorner u.
\]

为计算 \(X_{+}\) 和 \(X_{-}\)，以下述方式引入 \(\bigwedge V\) 的一个基：对于由 \([1, l]\) 的三个不相交子集组成的任意三元组 (A, B, C)，置

\[
e _ {\mathrm{A,B,C}} = e _ {a _ {1}} \wedge \dots \wedge e _ {a _ {m}} \wedge e _ {- b _ {1}} \wedge \dots \wedge e _ {- b _ {n}} \wedge e _ {c _ {1}} \wedge e _ {- c _ {1}} \wedge \dots \wedge e _ {c _ {p}} \wedge e _ {- c _ {p}}
\]

其中 \((a_{1},\ldots,a_{m})\)（分别为 \((b_{1},\ldots,b_{n}),(c_{1},\ldots,c_{p})\)）是 A（分别是 B、C）的元素按递增顺序排列所得的序列。这样得到 \(\bigwedge V\) 的一个基，简单计算表明

\[
X _ {-}. e _ {\mathrm{A,B,C}} = \sum_ {j \in \{1, l \}, j \notin \mathrm{A} \cup \mathrm{B} \cup \mathrm{C}} e _ {\mathrm{A,B,C} \cup \{j \}}\tag{8}
\]

\[
X _ {+}. e _ {\mathrm{A,B,C}} = - \sum_ {j \in \mathrm{C}} e _ {\mathrm{A,B,C-} \{j \}}.\tag{9}
\]

令 H 为 \(\bigwedge V\) 的自同态，它在 \((l - r)\) 上化为乘以 \(\bigwedge^{r}V\) 的乘法（\(0 \leq r \leq 2l\)）。容易验证（参见~习题19），有

\[
\begin{array}{c} {[ X _ {+}, X _ {-} ] = - H} \\ {[ H, X _ {+} ] = 2 X _ {+}} \\ {[ H, X _ {-} ] = - 2 X _ {-}.} \end{array}
\]

换言之，向量子空间 \(\mathfrak{s}\) 由 \(X_{+}, X_{-}\) 和 \(H\) 生成，且是 \(\operatorname{End}(\bigwedge \mathrm{V})\) 的李子代数，与 \(\mathfrak{sl}(2, k)\) 同构，而 \(\bigwedge^{r} \mathrm{V}\) 是权为 \(l - r\) 的元素子空间。令 \(\mathrm{E}_{r}\) 为 \(\bigwedge^{r} \mathrm{V}\) 中的本原元组成的子空间，即 \(\mathrm{E}_{r} = (\bigwedge^{r} \mathrm{V}) \cap \mathrm{Ker} X_{+}\)。由第1节可知，当 \(r < l\) 时，\(X_{-}\) 在 \(\bigwedge^{r} \mathrm{V}\) 上的限制是单射，并且当 \(r \leq l\) 时，\(\bigwedge^{r} \mathrm{V}\) 分解为直和。

\[
\begin{array}{c} \bigwedge^ {r} \mathrm{V} = \mathrm{E} _ {r} \oplus X _ {-} (\mathrm{E} _ {r - 2}) \oplus X _ {-} ^ {2} (\mathrm{E} _ {r - 4}) \oplus \dots \\ = \mathrm{E} _ {r} \oplus X _ {-} (\bigwedge^ {r - 2} \mathrm{V}). \end{array}
\]

这尤其表明，对于 \(\dim \mathrm{E}_r = \binom{2l}{r} - \binom{2l}{r-2}\)，\(0 \leq r \leq l\)。

另一方面，\(\mathfrak{sp}(\varPsi)\) 的定义本身表明，\(\Gamma^{*}\) 被 \(\sigma\) 的对偶的第二外幂湮灭。同样，\(\Gamma\) 被 \(\bigwedge^{2}\sigma\) 湮灭。立即得到，\(X_{+}\) 和 \(X_{-}\)，从而 \(H\)，也与 \(\bigwedge\sigma(g)\) 所对应的自同态 \(g\in\mathfrak{g}\) 交换。因此，对于 \(\mathrm{E}_{r}\)，子空间 \(0\leq r\leq l\) 在 \(\bigwedge^{r}\sigma\) 下稳定；我们将证明，\(\bigwedge^{r}\sigma\) 在 \(\mathrm{E}_{r}\) 上的限制是权为 \(\sigma_{r}\) 的基本表示 \(\varpi_{r}\)（\(1\leq r\leq l\)）。

我们先注意，\(\bigwedge^{r}\sigma\) 相对于 \(\mathfrak{h}\) 的权是

\[
\varepsilon_ {i _ {1}} + \dots + \varepsilon_ {i _ {k}} - (\varepsilon_ {j _ {1}} + \dots + \varepsilon_ {j _ {r - k}}),
\]

其中 \(i_1, \ldots, i_k\)（分别为 \(j_1, \ldots, j_{r-k}\)）是 \([1, l]\) 中的互异元素；因此，\(\bigwedge^r \sigma\) 的最高权确实是

\[
\varpi_ {r} = \varepsilon_ {1} + \dots + \varepsilon_ {r}
\]

表示 \(\varpi_{r}\) 的最高权向量是 \(e_{1} \wedge \cdots \wedge e_{r} = e_{\{1,\ldots,r\},\varnothing,\varnothing}\)，且它属于 \(e_{1} \wedge \cdots \wedge e_{r} \in E_{r}\)。由于 \(\bigwedge^{r}\sigma\) 是不可约表示，\(E_{r}\) 是其最高权

若 \(s \in \mathbf{Sp}(\varPsi)\)，则 s 在 \(\bigwedge V\)（分别在 \(\bigwedge V^{*}\)）上的延拓固定 \(\varGamma\)（分别固定 \(\varGamma^{*}\)），因而与 \(X_{+}\) 和 \(X_{-}\) 交换，并使 \(E_{r}\) 稳定。因此，\(E_{r}\) 包含由 \(F_{r}\) 将 \(e_{1} \wedge \cdots \wedge e_{r}\) 变换所得的向量生成的子空间 \(\mathbf{Sp}(\varPsi)\)。Witt 定理表明，这些变换得到的正是非零可分解 r-向量，且其对应的 V 的向量子空间是全迷向子空间；我们称这样的 r-向量为迷向的。

引理 2. 对于 \(1 \leq r \leq l\)，令 \(\mathrm{F}_r\) 为由迷向 \(\bigwedge^r\mathrm{V}\)-向量生成的 \(r\) 的子空间。则

\[
\bigwedge^ {r} \mathrm{V} = \mathrm{F} _ {r} + X _ {-} (\bigwedge^ {r - 2} \mathrm{V}) = \mathrm{F} _ {r} + \left(\sum_ {i = 1} ^ {l} e _ {i} \wedge e _ {- i}\right) \wedge \bigwedge^ {r - 2} \mathrm{V}.
\]

由于 \(F_{r} \subset E_{r}\)，且 \(\mathrm{E}_{r} \cap X_{-}(\bigwedge^{r-2}\mathrm{V}) = \{0\}\)，引理推出 \(F_{r} = E_{r}\)。另一方面，令 \(s \in \mathbf{Sp}(\varPsi)\)；自同构 \(a \mapsto sas^{-1}\) 将 \(\operatorname{End}(V)\) 保持为自身，保持 g 不变，并在其上诱导 \(\operatorname{Aut}_{0}(\mathfrak{g})\) 的一个元素（参见~(VII)），因而将 g 的每个不可约表示变为等价表示（第7节，第1节，命题2）。由于 \(e_{1} \wedge \cdots \wedge e_{r}\) 属于 \(\bigwedge^{r}\sigma\) 的一个不可约分支，且 \(E_{r} = F_{r}\) 由 \(e_{1} \wedge \cdots \wedge e_{r}\) 对 \(\mathbf{Sp}(\varPsi)\) 的变换生成，所以 g 在 \(E_{r}\) 上的表示是同构型的。但是，其最高权 \(\varpi_{r}\) 的重数为 1，故它不可约。

引理尚待证明。r = 1 时显然。作归纳并设 \(r \geq 2\)。根据归纳假设，只需证明

\[
\mathrm{F} _ {r - 1} \wedge \mathrm{V} \subset \mathrm{F} _ {r} + \Gamma \wedge \bigwedge^ {r - 2} \mathrm{V},
\]

或者说，若 \(y\) 是一个可分解的 \((r - 1)\) -向量且 \(x \in V\)，则

\[
z = y \wedge x \in \mathrm{F} _ {r} + \Gamma \wedge \bigwedge^ {r - 2} \mathrm{V}.
\]

令 \((f_{i})_{1\leq\pm i\leq l}\) 为 V 的 Witt 基，使 \(y=f_{1}\wedge\cdots\wedge f_{r-1}\)。只需在 \(x=f_{i}\) 时完成证明。若 \(i\notin(1-r,-1)\)，r-向量 \(f_{1}\wedge\cdots\wedge f_{r-1}\wedge f_{i}\) 是迷向的。否则，必要时重新编号 \(f_{i}\)，可设 i=1-r。于是 \(\Gamma=\sum_{j=1}^{l}f_{j}\wedge f_{-j}\)，所以

\[
\begin{array}{l} f _ {r - 1} \wedge f _ {1 - r} = \frac {1}{l - r + 2} \left(\Gamma - \sum_ {i = 1} ^ {r - 2} f _ {i} \wedge f _ {- i} + \sum_ {j = r} ^ {l} (f _ {r - 1} \wedge f _ {1 - r} - f _ {j} \wedge f _ {- j})\right), \\ z = \frac {1}{l - r + 2} \Gamma \wedge f _ {1} \wedge \dots \wedge f _ {r - 2} \\ \qquad + \frac {1}{l - r + 2} \sum_ {j = r} ^ {l} (f _ {1} \wedge \dots \wedge f _ {r - 2}) \wedge (f _ {r - 1} \wedge f _ {1 - r} - f _ {j} \wedge f _ {- j}). \end{array}
\]

但是

\[
f _ {r - 1} \wedge f _ {1 - r} - f _ {j} \wedge f _ {- j} = (f _ {r - 1} + f _ {j}) \wedge (f _ {1 - r} - f _ {- j}) - f _ {j} \wedge f _ {1 - r} + f _ {r - 1} \wedge f _ {- j}
\]

并立即验证下列 r-向量

\[
\begin{array}{c} f _ {1} \wedge \dots \wedge f _ {r - 2} \wedge (f _ {r - 1} + f _ {j}) \wedge (f _ {1 - r} - f _ {- j}), \\ f _ {1} \wedge \dots \wedge f _ {r - 2} \wedge f _ {j} \wedge f _ {1 - r} \text {and} f _ {1} \wedge \dots \wedge f _ {r - 1} \wedge f _ {- j} \end{array}
\]

当 \(r \leq j \leq l\) 时是迷向的。因此，\(z \in F_{r} + \Gamma \wedge \bigwedge^{r-2} V\)，从而证明完毕。

\begin{enumerate}
\def\labelenumi{(\Alph{enumi})}
\setcounter{enumi}{21}
\tightlist
\item
有 \(w_{0} = -1\)，所以 g 的每个有限维单表示都是正交型或辛型。由第 VI 章，§4，第6.VI节，\(\varpi_{r}\) 关于 \((\alpha_{1}, \ldots, \alpha_{l})\) 的坐标之和为
\end{enumerate}

\[
1 + 2 + \dots + (r - 1) + r + r + \dots + r + \frac {r}{2}
\]

因此 \(\sigma_{r}\) 在 r 为偶数时是正交型的，在 r 为奇数时是辛型的。

由于 \(e_{1} \wedge \cdots \wedge e_{r}\) 和 \(e_{-1} \wedge \cdots \wedge e_{-r}\) 属于 \(E_{r}\)，且

\[
\Psi_ {(r)} \left(e _ {1} \wedge \dots \wedge e _ {r}, e _ {- 1} \wedge \dots \wedge e _ {- r}\right) = 1,
\]

可知 \(\Psi_{(r)}\) 在 \(E_{r}\) 上的限制非零：这在相差一个非零常数因子的意义下就是该双线性型；当 r 为偶数时它对称，当 r 为奇数时它交错，并且在 \(\sigma_{r}\) 下不变。

\begin{enumerate}
\def\labelenumi{(\Roman{enumi})}
\setcounter{enumi}{5}
\tightlist
\item
对于所有 \(x \in \mathfrak{g}\)，\(\sigma(x)\) 的特征多项式具有形式
\end{enumerate}

\[
\mathrm{T} ^ {2 l} + f _ {1} (x) \mathrm{T} ^ {2 l - 1} + \dots + f _ {2 l} (x)
\]

其中 \(f_{1},\ldots ,f_{2l}\) 是 \(\mathfrak{g}\) 上的不变多项式函数。

若 \(x = \xi_{1}H_{1} + \cdots + \xi_{l}H_{l} \in \mathfrak{h}\)，则 \(f_{i}(x)\) 在相差一个符号的意义下是 \(\xi_{1}, \ldots, \xi_{l}, -\xi_{1}, \ldots, -\xi_{l}\) 的初等对称函数；这些对称函数在奇数次数时为零，并且

\[
\mathrm{T} ^ {2 l} + f _ {2} (x) \mathrm{T} ^ {2 l - 2} + \dots + f _ {2 l} (x) = (\mathrm{T} ^ {2} - \xi_ {1} ^ {2}) \ldots (\mathrm{T} ^ {2} - \xi_ {l} ^ {2}).
\]

像第2.VI节中那样，可知 \(f_{1}=f_{3}=f_{5}=\cdots=0\)，并且

\[
(f _ {2}, f _ {4}, \dots , f _ {2 l})
\]

是生成 \(\mathfrak{g}\) 上不变多项式函数代数的代数自由族。

\begin{enumerate}
\def\labelenumi{(\Roman{enumi})}
\setcounter{enumi}{6}
\tightlist
\item
由于邓金图的唯一自同构是恒等自同构，有 \(\mathrm{Aut}(\mathfrak{g}) = \mathrm{Aut}_0(\mathfrak{g})\)。
\end{enumerate}

令 \(\Sigma\) 为 \(V\) 关于 \(\Psi\) 的相似变换群。如第 2.VII 节中那样可证明，\(\mathfrak{g}\) 的自同构是映射 \(x \mapsto sxs^{-1}\)，其中 \(s \in \Sigma\)，所以 \(\operatorname{Aut}(\mathfrak{g}) = \operatorname{Aut}_0(\mathfrak{g})\) 可与 \(\Sigma / k^*\) 识别。

对于所有 \(s \in \Sigma\)，令 \(\mu(s)\) 为 s 的乘子。映射 \(s \mapsto \mu(s) \mod k^{*2}\) 从 \(\Sigma\) 到 \(k^{*}/k^{*2}\) 是同态，其核包含 \(k^{*}.1\)，因而给出从 \(\lambda\) 到 \(\Sigma/k^{*}\) 的同态 \(k^{*}/k^{*2}\)。有 \(\mathbf{Sp}(\varPsi) \cap k^{*} = \{1, -1\}\)。考虑同态序列

\[
1 \longrightarrow \mathbf {S p} (\Psi) / \{1, - 1 \} \stackrel {{\iota}} {{\longrightarrow}} \Sigma / k ^ {*} \stackrel {{\lambda}} {{\longrightarrow}} k ^ {*} / k ^ {* 2} \longrightarrow 1.\tag{10}
\]

映射 \(\iota\) 是单射，且 \(\mathrm{Im}(\iota) \subset \mathrm{Ker}\lambda\) 包含于 \(\mathbf{Sp}(\varPsi)\)，其元素的乘子为 1。如果 \(s \in \varSigma\) 的乘子是 \(k^{*2}\) 中的元素，则存在 \(\nu \in k^*\)，使 \(\nu s \in \mathbf{Sp}(\varPsi)\)；因此 \(\mathrm{Im}(\iota) = \mathrm{Ker}(\lambda)\)。总之，序列（10）是正合的。将 \(\mathbf{Sp}(\varPsi)/\{1, -1\}\) 识别为 \(\varSigma/k^*\) 的一个子群。由于 \(k^*/k^{*2}\) 是交换群，\(\mathbf{Sp}(\varPsi)/\{1, -1\}\) 包含 \(\varSigma/k^*\) 的换位子群。因此，\(\mathrm{Aut}_e(\mathfrak{g})\) 包含于 \(\mathbf{Sp}(\varPsi)/\{1, -1\}\)（\(\S\) 11，第2节，命题3）。事实上，它等于该群，并且 \(\mathrm{Aut}(\mathfrak{g})/\mathrm{Aut}_e(\mathfrak{g})\) 可与 \(k^*/k^{*2}\) 识别（习题9）。

\begin{enumerate}
\def\labelenumi{(\Roman{enumi})}
\setcounter{enumi}{7}
\tightlist
\item
典则双线性型 \(\Phi_{\mathrm{R}}\) 在 \(\mathfrak{h}^*\) 上为
\end{enumerate}

\[
\varPhi_ {\mathrm{R}} (\xi_ {1} \varepsilon_ {1} + \dots + \xi_ {l} \varepsilon_ {l}, \xi_ {1} ^ {\prime} \varepsilon_ {1} + \dots + \xi_ {l} ^ {\prime} \varepsilon_ {l}) = \frac {1}{4 (l + 1)} (\xi_ {1} \xi_ {1} ^ {\prime} + \dots + \xi_ {l} \xi_ {l} ^ {\prime})
\]

（第 VI 章，§4，第6.V节）。因此，\(\varPhi_{\mathrm{R}}\) 的逆型，即 Killing 型在 \(\mathfrak{h}\) 上的限制，为

\[
\Phi (\xi_ {1} H _ {1} + \dots + \xi_ {l} H _ {l}, \xi_ {1} ^ {\prime} H _ {1} + \dots + \xi_ {l} ^ {\prime} H _ {l}) = 4 (l + 1) (\xi_ {1} \xi_ {1} ^ {\prime} + \dots + \xi_ {l} \xi_ {l} ^ {\prime}).
\]

\begin{enumerate}
\def\labelenumi{(\Roman{enumi})}
\setcounter{enumi}{8}
\tightlist
\item
回顾由公式（6）定义的 \(X_{\alpha}\)（\(\alpha \in R\)）。容易验证，对于 \([X_{\alpha}, X_{-\alpha}] = -H_{\alpha}\)，有 \(\alpha \in R\)。另一方面，映射 \(\theta : a \mapsto -^{t}a\) 是 g 的自同构，且对所有 \(\theta(X_{\alpha}) = X_{-\alpha}\) 有 \(\alpha \in R\)。因此，\((X_{\alpha})_{\alpha \in \mathbb{R}}\) 是 \((\mathfrak{g}, \mathfrak{h})\) 中的谢瓦莱系。
\end{enumerate}

假设 \(k = \mathbf{Q}\)。嘉当子代数 \(\mathfrak{h}\) 有两个许可格：\(\mathrm{Q}(\mathrm{R}^{\vee}) = \sum_{i=1}^{l} \mathbf{Z}.H_i\) 和 \(\mathrm{P}(\mathrm{R}^{\vee}) = \mathrm{Q}(\mathrm{R}^{\vee}) + \frac{1}{2}\mathbf{Z}. \sum_{i=1}^{l} H_i\)。可见，\(\mathrm{Q}(\mathrm{R}^{\vee})\) 是 \(\mathfrak{h}\) 中具有整数项的矩阵集合。因此，谢瓦莱序 \(\mathrm{Q}(\mathrm{R}^{\vee}) + \sum_{\alpha \in \mathrm{R}} \mathbf{Z}.X_{\alpha}\) 是 \(\mathfrak{sp}(2l,\mathbf{Z})\)（其中矩阵属于 \(\mathfrak{g}\)）。

考虑约化李代数 \(\mathfrak{sp}(\varPsi) + \mathbf{Q}.1\)。容易看出，其具有整数项的元素组成一个谢瓦莱序；该序平行于 \(\mathfrak{sp}(\varPsi)\) 投影到 \(\mathbf{Q}.1\) 上，所得正是谢瓦莱序 \(\mathrm{P}(\mathrm{R}^{\vee}) + \sum \mathbf{Z}.X_{\alpha}\)。

最后，\(X_{\alpha}^{2} = 0\) 对所有 \(\alpha \in \mathbb{R}\) 都有。由此，格 \(\overline{\mathcal{V}}\) 位于 \(\mathrm{V}\) 中，并由 Witt 基 \(e_i\) 生成；它是 \(\mathfrak{sp}(2l,\mathbf{Z})\) 的容许格。在 \(\mathrm{E}_r \cap \bigwedge^r \mathcal{V}\) 中的 \(\mathrm{E}_r\) 也如此。

最后，\(E_{r}\) 对谢瓦莱序

\[
\mathrm{P} (\mathrm{R} ^ {\vee}) + \sum \mathbf {Z}. X _ {\alpha}
\]

仅当 \(r\) 为偶数时才有；此时 \(\mathrm{E}_r\cap \bigwedge^r\mathcal{V}\) 就是这样的格。

\subsection{\texorpdfstring{类型为 \(D_{l}\) 的代数 \((l\geq 2)\)}{类型为 D\_\{l\} 的代数 (l\textbackslash geq 2)}}

\begin{enumerate}
\def\labelenumi{(\Roman{enumi})}
\tightlist
\item
  设 V 为维数为 \(2l \geq 4\) 的向量空间，且 \(\varPsi\) 为 V 上指标为最大指标 l 的非退化对称双线性型。由《代数》第 IX 章，§4，第2节，V 可写成两个极大全迷向子空间 F 和 \(F'\) 的直和。令 \((e_i)_{1 \leq i \leq l}\) 为 F 的一个基，令 \((e_{-i})_{1 \leq i \leq l}\) 为 \(F'\) 的对偶基（关于由 \(F'\) 定义的 F 与 \(\varPsi\) 之间的对偶关系）。于是 \(e_1, \ldots, e_l, e_{-l}, \ldots, e_{-1}\) 是 V 的一个基；我们称其为 V 的 Witt 基。相对于此基，\(\varPsi\) 的矩阵是阶为 2l 的方阵 S，其元素全为 0，唯有位于第二对角线上的元素等于 1。代数 \(\mathfrak{g} = \mathfrak{o}(\varPsi)\) 可与方阵代数 \(\mathfrak{o}_S(2l, k)\) 识别，后者由满足 \(g = -S^t g S\) 的阶为 2l 的方阵 g 组成。其维数为 \(l(2l - 1)\)。直接计算表明，g 是如下形式的矩阵的集合
\end{enumerate}

\[
\left( \begin{array}{c c} A & B \\ C & D \end{array} \right)
\]

其中 A、B、C、D 是阶为 l 的方阵，并满足 \(B = -s^{t}Bs\)、\(C = -s^{t}Cs\) 和 \(D = -s^{t}As\)（s 是阶为 l 的矩阵，除位于第二对角线上的元素等于 1 外，其余元素全为 0）。

令 \(\mathfrak{h}\) 为属于 \(\mathfrak{g}\) 的对角矩阵的集合。这是 \(\mathfrak{g}\) 的交换子代数，其基由元素 \(H_{i} = E_{i,i} - E_{-i,-i}\)（\(1 \leq i \leq l\)）组成。令 \((\varepsilon_{i})\) 为与 \(\mathfrak{h}^{*}\) 对偶的 \((H_{i})\) 的基。对于 \(1 \leq i < j \leq l\)，置

\[
\left\{ \begin{array}{l l} X _ {\varepsilon_ {i} - \varepsilon_ {j}} & = E _ {i, j} - E _ {- j, - i} \\ X _ {- \varepsilon_ {i} + \varepsilon_ {j}} & = - E _ {j, i} + E _ {- i, - j} \\ X _ {\varepsilon_ {i} + \varepsilon_ {j}} & = E _ {i, - j} - E _ {j, - i} \\ X _ {- \varepsilon_ {i} - \varepsilon_ {j}} & = - E _ {- j, i} + E _ {- i, j}. \end{array} \right.\tag{11}
\]

这些元素构成 h 在 g 中的一个补空间的基。对于 \(h \in h\)，有

\[
[ h, X _ {\alpha} ] = \alpha (h) X _ {\alpha}
\]

对所有 \(\alpha \in \mathbb{R}\) 都成立，其中 \(\mathbb{R}\) 是由 \(\pm \varepsilon_i \pm \varepsilon_j\)（\(i < j\)）组成的集合。因此，\(\mathfrak{h}\) 是 \(\mathfrak{g}\) 的分裂嘉当子代数，且 \((\mathfrak{g}, \mathfrak{h})\) 的根是 \(\mathbb{R}\) 的元素。\(\mathbb{R}\) 的根系 \((\mathfrak{g}, \mathfrak{h})\) 在 \(\mathrm{D}_l\) 时为 \(l \geq 3\) 型，在 \(\mathrm{A}_1 \times \mathrm{A}_1\) 时为 \(\mathrm{D}_2\) 型（即 \(l = 2\) 型）（第 VI 章，§4，第8.I节，推广到 \(l = 2\) 的情形）。因此，\(\mathfrak{g}\) 是 \(\mathrm{D}_l\) 型的可分裂单李代数（当 \(l \geq 3\) 时）。

g 的每个分裂嘉当子代数都可由 g 的一个初等自同构变换为 h，因而也可由 \(\mathbf{O}(\varPsi)\) 的某个元素变换得到（参见~(VII)）；所以它是集合 \(h_{\beta}\)，其中的元素 g 相对于 V 的某个 Witt 基 \(\beta\) 的矩阵为对角矩阵。立即验证可知，在 \(h_{\beta}\) 下不变的唯一子空间是由 \(\beta\) 的某个子集生成的子空间。

由于代数 \(\mathfrak{o}_{S}(4,k)\) 和 \(\mathfrak{sl}(2,k)\times\mathfrak{sl}(2,k)\) 具有相同的根系，二者同构。同样，\(\mathfrak{o}_{S}(6,k)\) 和 \(\mathfrak{sl}(4,k)\) 同构（另参见~习题3）。下文设 \(l\geq3\)。

\begin{enumerate}
\def\labelenumi{(\Roman{enumi})}
\setcounter{enumi}{1}
\tightlist
\item
我们通过第 VI 章，§4，第8.V节确定 \(\mathbb{R}^{\vee}\)。我们得到
\end{enumerate}

\[
H _ {\varepsilon_ {i} - \varepsilon_ {j}} = H _ {i} - H _ {j}, \quad H _ {\varepsilon_ {i} + \varepsilon_ {j}} = H _ {i} + H _ {j}.
\]

\begin{enumerate}
\def\labelenumi{(\Roman{enumi})}
\setcounter{enumi}{2}
\tightlist
\item
置 \(\alpha_{1} = \varepsilon_{1} - \varepsilon_{2},\alpha_{2} = \varepsilon_{2} - \varepsilon_{3},\ldots ,\alpha_{l - 1} = \varepsilon_{l - 1} - \varepsilon_{l},\alpha_{l} = \varepsilon_{l - 1} + \varepsilon_{l}\)。由第 VI 章，§4，第8.II节，\((\alpha_{1},\dots,\alpha_{l})\) 是 R 的一个基 B；相对于 B 的正根是 \(\varepsilon_i\pm \varepsilon_j\)（\(i < j\)）。相应的 Borel 子代数 \(\mathfrak{b}\) 是属于 \(\mathfrak{g}\) 的上三角矩阵的集合。
\end{enumerate}

立即验证可知，在 b 下不变的非平凡向量子空间只有全迷向子空间 \(V_{1},\ldots,V_{l},V_{l}^{\prime}\)，其中 \(V_{i}\) 由 \(e_{1},\ldots,e_{i}\) 生成，\(V_{l}^{\prime}\) 由 \(e_{1},\ldots,e_{l-1},e_{-l}\) 生成，以及 \(V_{-1},\ldots,V_{-l+1}\) 的正交补 \(V_{1},\ldots,V_{l-1}\)；\(V_{-i}\) 的正交补 \(V_{i}\) 由 \(e_{1},\ldots,e_{l},e_{-l},\ldots,e_{-(i+1)}\) 生成。但是直接计算表明，若 \(a\in g\) 使 \(V_{l-1}\) 稳定，则其矩阵具有形式

\[
\left( \begin{array}{c c c} A & x & B \\ 0 & \left( \begin{array}{c c} \lambda & 0 \\ 0 & - \lambda \end{array} \right) & y \\ 0 & 0 & D \end{array} \right)
\]

其中 \(A, B, D\) 是阶为 \(l - 1\) 的方阵，\(x\)（分别 \(y\)）是具有 2 列和 \(l - 1\) 行（分别具有 2 行和 \(l - 1\) 列）的矩阵，且 \(\lambda \in k\)。由此 \(a\) 使 \(\mathrm{V}_l\) 和 \(\mathrm{V}_l'\) 稳定。因此，\(\mathfrak{b}\) 是 \(a \in \mathfrak{g}\) 中使迷向旗 \((\mathrm{V}_1, \ldots, \mathrm{V}_{l-1})\) 的所有元素稳定的集合。注意，前述事实和 Witt 定理（...）表明，唯一的相关极大全迷向子空间是 \(\mathrm{V}_l\) 和 \(\mathrm{V}_l'\)，均包含 \(\mathrm{V}_{l-1}\)。

我们称一个迷向旗为拟极大迷向旗，如果它由 l-1 个维数为 \(1,\ldots,l-1\) 的全迷向子空间组成。于是，如第2节中那样可知，对于任意拟极大迷向旗 \(\delta\)，使 \(b_{\delta}\) 的各元素稳定的 \(a\in g\) 的集合 \(\delta\) 是 g 的一个 Borel 子代数，并且映射 \(\delta\mapsto b_{\delta}\) 是从拟极大迷向旗集合到 Borel 子代数集合的双射。

我们称一个迷向旗为真的，如果它不同时含有一个 \(l\) 维子空间和一个 \(l-1\) 维子空间。令 \(\delta\) 为这样的迷向旗，并令 \(\mathfrak{p}_{\delta}\) 为由 \(a \in \mathfrak{g}\) 中使 \(\delta\) 的各元素稳定的元素组成的集合。若它包含于 \(\delta \subset \{\mathrm{V}_1, \ldots, \mathrm{V}_l, \mathrm{V}_l'\}\)，则 \(\mathfrak{p}_{\delta}\) 是 \(\mathfrak{g}\) 的一个抛物子代数，包含 \(\mathfrak{b}\)，并且容易验证，唯一不等于 \(\neq \{0\}\) 且在 \(\mathfrak{p}_{\delta}\) 下稳定的全迷向子空间是 \(\delta\) 的元素。由于每类有 \(2^{l-2}\) 个 \(\{\mathrm{V}_1, \ldots, \mathrm{V}_l, \mathrm{V}_l'\}\) 中的真的迷向旗，其中包含 \(\mathrm{V}_{l-1}\)（分别包含 \(\mathrm{V}_l\)、包含 \(\mathrm{V}_l'\)、既不包含 \(\mathrm{V}_{l-1}\)、也不包含 \(\mathrm{V}_l\) 和 \(\mathrm{V}_l'\)），这给出了 \(2^l\) 个抛物子代数。由此，如上所述，映射 \(\mathfrak{b}\) 是从真的迷向旗集合到 \(\delta \mapsto \mathfrak{p}_{\delta}\) 的抛物子代数集合（属于 \(\mathfrak{g}\)）的双射。

\begin{enumerate}
\def\labelenumi{(\Roman{enumi})}
\setcounter{enumi}{3}
\tightlist
\item
与 \(\alpha_{1},\ldots,\alpha_{l}\) 对应的基本权由第 VI 章，§4，第8.VI节给出：
\end{enumerate}

\[
\varpi_ {i} = \varepsilon_ {1} + \varepsilon_ {2} + \dots + \varepsilon_ {i} (1 \leq i \leq l - 2)
\]

\[
\begin{array}{c} \varpi_ {l - 1} = \frac {1}{2} (\varepsilon_ {1} + \varepsilon_ {2} + \ldots + \varepsilon_ {l - 2} + \varepsilon_ {l - 1} - \varepsilon_ {l}) \\ \varpi_ {l} = \frac {1}{2} (\varepsilon_ {1} + \varepsilon_ {2} + \ldots + \varepsilon_ {l - 2} + \varepsilon_ {l - 1} + \varepsilon_ {l}). \end{array}
\]

令 \(\sigma\) 为 \(\mathfrak{g}\) 在 V 上的恒等表示。外幂 \(\bigwedge^{r}\sigma\) 作用于 \(\mathrm{E} = \bigwedge^{r}(\mathrm{V})\)。若 \(h\in \mathfrak{h}\)，有

\[
\sigma (h) e _ {i} = \varepsilon_ {i} (h) e _ {i}, \quad \sigma (h) e _ {- i} = - \varepsilon_ {i} (h) e _ {- i}
\]

对于 \(1 \leq i \leq l\)，由此，对于 \(1 \leq r \leq l\)，\(\varepsilon_{1} + \cdots + \varepsilon_{r}\) 是 \(\bigwedge^{r} \sigma\) 的最高权，权为 \(\varepsilon_{1} + \cdots + \varepsilon_{r}\) 的元素正是与 \(e_{1} \wedge \cdots \wedge e_{r}\) 成比例的元素。

我们将证明，当 \(1 \leq r \leq l - 2\) 时，表示 \(\bigwedge^{r} \sigma\) 是权为 \(\varpi_{r}\) 的基本表示。

为此，只需证明 \(\bigwedge^{r}\sigma\) 在 \(1 \leq r \leq l - 1\) 时并非不可约（注意表示 \(\bigwedge^{l}\sigma\) 不可约，参见~习题10），或者证明 \(T_{r}\) 中包含 \(\bigwedge^{r}V\) 且在 g 下稳定的最小子空间 \(e_{1} \wedge \cdots \wedge e_{r}\) 就是整个 \(\bigwedge^{r}V\)。当 r = 1 时这立即成立。当 r = 2 时，如第2节中那样，我们看到 \(\bigwedge^{2}\sigma\) 等价于 g 的伴随表示；由于 g 是单的，该表示不可约。通过像第2节中那样对 l 作归纳，并设 \(l - 1 \geq r \geq 3\)，即可完成证明。

现在确定最高权为 \(\varpi_{l-1}\) 和 \(\varpi_{l}\) 的基本表示。令 Q 为二次型 \(x \mapsto \frac{1}{2}\Psi(x, x)\)。我们在第2.IV节中定义了 Clifford 代数 \(\lambda\) 在 \(\mathrm{C}(\mathrm{Q})\) 上的旋量表示 \(N = \bigwedge F'\)。立即验证可知，N 的子空间 \(N_{+}\)（分别为 \(N_{-}\)）由 p 为偶数（分别为奇数）时的 \(\bigwedge^{p}F'\) 之和给出，并且在 \(\lambda\) 对 \(\mathrm{C}^{+}(\mathrm{Q})\) 的限制下稳定。因此，\(\lambda_{+}\) 和 \(\lambda_{-}\) 分别是 \(\mathrm{C}^{+}(\mathrm{Q})\) 在 \(N_{+}\) 和 \(N_{-}\) 上的表示，也是 \(\mathrm{C}^{+}(\mathrm{Q})\) 的半旋量表示（《代数》第 IX 章，§9，第4节）；它们不可约，维数为 \(2^{l-1}\)，且不等价。令 \(\rho_{+} = \lambda_{+} \circ f\) 和 \(\rho_{-} = \lambda_{-} \circ f\) 为相应的 g 的不可约表示（第2节，引理1 (vi)）。根据引理1 (i)，有

\[
f (H _ {i}) = \frac {1}{2} (e _ {i} e _ {- i} - e _ {- i} e _ {i}) = e _ {i} e _ {- i} - \frac {1}{2} = \frac {1}{2} - e _ {- i} e _ {i}
\]

并且如第2.IV节中那样可见，对于 \(h \in h\) 和 \(1 \leq i_{1} < \cdots < i_{k} \leq l\)，

\[
\begin{array}{l} \lambda \circ f (h) (e _ {- i _ {1}} \wedge \dots \wedge e _ {- i _ {k}}) \\ = (\frac {1}{2} (\varepsilon_ {1} + \dots + \varepsilon_ {l}) - (\varepsilon_ {i _ {1}} + \dots + \varepsilon_ {i _ {k}})) (h) (e _ {- i _ {1}} \wedge \dots \wedge e _ {- i _ {k}}). \end{array}
\]

因此，\(\rho_{+}\)（分别 \(\rho_{-}\)）的最高权是 \(\varpi_{l}\)（分别 \(\varpi_{l-1}\)）。

我们称 \(\rho_{+}\) 和 \(\rho_{-}\) 为 g 的半旋量表示。它们的所有权都是单权。我们还称 \(\rho = \lambda \circ f = \rho_{+} \oplus \rho_{-}\) 为 g 的旋量表示。

\begin{enumerate}
\def\labelenumi{(\Alph{enumi})}
\setcounter{enumi}{21}
\tightlist
\item
对于 \(1 \leq r \leq l - 2\)，基本表示 \(\bigwedge^{r} \sigma\) 是正交型的：它保持 \(\varPsi\) 在 \(\bigwedge^{r} \mathrm{V}\) 上的延拓不变。
\end{enumerate}

现在考虑 \(\rho\) 的旋量表示 \(\mathfrak{g}\)。如第2节中那样，我们说明，对于 \(1 \leq i, j \leq l, i \neq j\)，

\[
\begin{array}{r} f (X _ {\varepsilon_ {i} - \varepsilon_ {j}}) = \pm e _ {i} e _ {- j} \\ f (X _ {\varepsilon_ {i} + \varepsilon_ {j}}) = \pm e _ {i} e _ {j} \\ f (X _ {- \varepsilon_ {i} - \varepsilon_ {j}}) = \pm e _ {- i} e _ {- j}. \end{array}
\]

由此，在第2.V节中引入的非退化双线性型 \(\varPhi\) 在 \(\rho(\mathfrak{g})\) 下不变。因此，旋量表示 \(\rho\) 保持一个非退化型不变；当 \(l \equiv 0, -1 \pmod{4}\) 时该型对称，当 \(l \equiv 1, 2 \pmod{4}\) 时该型交错。

当 \(l\) 为奇数时，\(\varPhi\) 交换 \(\mathrm{N}_{+}\) 与 \(\mathrm{N}_{-}\)；当 \(l \equiv 0 \pmod{4}\) 时，它具有相应的性质；当 \(l \equiv 2 \pmod{4}\) 时，\(w_{0} = -1\) 具有所述作用。

另一方面，若 \(l\) 为奇数，则 \(\mathrm{N}_{+}\) 和 \(\mathrm{N}_{-}\) 关于 \(\varPhi\) 都具有相应的性质。此外，\(-w_{0}(\alpha_{i}) = \alpha_{i}\) 对 \(1 \leq i \leq l - 2\) 成立，\(-w_{0}(\alpha_{l}) = \alpha_{l - 1}\)，而 \(-w_{0}(\alpha_{l - 1}) = \alpha_{l}\)，所以 \(-w_{0}(\varpi_{l}) = \varpi_{l - 1}\)。

\begin{enumerate}
\def\labelenumi{(\Roman{enumi})}
\setcounter{enumi}{5}
\tightlist
\item
对于所有 \(x \in \mathfrak{g}\)，\(\sigma(x)\) 的特征多项式具有形式
\end{enumerate}

\[
\mathrm{T} ^ {2 l} + f _ {1} (x) \mathrm{T} ^ {2 l - 1} + \dots + f _ {2 l} (x).
\]

我们像第3.VI节中那样看到 \(f_{1}=f_{3}=f_{5}=\cdots=0\)。由第 VI 章，§4，第8.IX节和§8，第3节，定理1，存在 g 上的多项式函数 \(\tilde{f}\)，使得 \(f_{2}, f_{4}, \ldots, f_{2l-2}, \tilde{f}\) 生成 \(\mathrm{I}(\mathfrak{g}^{*})\)（即 g 上不变多项式函数的代数），它们代数无关，并且还有 \(\tilde{f}^{2}=(-1)^{l}f_{2l}\)。

对于 \(x \in \mathfrak{g}\)，有 \({}^t(Sx) = {}^t xS = -Sx\)，所以可以考虑 \(\operatorname{Pf}(Sx)\)，它是 \(x\) 的多项式函数。现在

\[
f _ {2 l} (x) = \det (x) = (- 1) ^ {l} \det (S x) = (- 1) ^ {l} (\operatorname{Pf} (S x)) ^ {2}.
\]

因此，可取 \(\tilde{f}(x) = \mathrm{Pf}(Sx)\)。

\begin{enumerate}
\def\labelenumi{(\Roman{enumi})}
\setcounter{enumi}{6}
\tightlist
\item
回顾（§5，第3节，命题5的推论1），\(\mathrm{Aut}(\mathfrak{g}) / \mathrm{Aut}_0(\mathfrak{g})\) 可与 \(\mathrm{Aut}(\mathrm{D})\) 识别，后者是 \((\mathfrak{g},\mathfrak{h})\) 的邓金图 D 的自同构群。当 \(l\neq 4\) 时，\(\mathrm{Aut}(\mathrm{D})\) 是阶为 2 的群，由\(\alpha_{1},\ldots ,\alpha_{l}\)的置换组成，这些置换保持\(\alpha_{1},\ldots ,\alpha_{l - 2}\)不变。当 \(l = 4\) 时，\(\mathrm{Aut}(\mathrm{D})\) 由保持 \(\alpha_{1},\ldots ,\alpha_{4}\) 不变的 \(\alpha_{2}\) 的置换组成；它同构于 \(\mathfrak{S}_3\)（参见~第 VI 章，§4，第8.XI节）。在所有情形下，\(\mathrm{Aut}(\mathrm{D})\) 中保持 \(\alpha_{1}\) 不变的元素组成的子群的阶为 2。我们将相应的 \(\operatorname{Aut}'(\mathfrak{g})\) 的子群记为 \(\operatorname{Aut}(\mathfrak{g})\)；当 \(\operatorname{Aut}'(\mathfrak{g}) = \operatorname{Aut}(\mathfrak{g})\) 时，\(l \neq 4\)，当 l = 4 时 \((\operatorname{Aut}(\mathfrak{g}) : \operatorname{Aut}'(\mathfrak{g})) = 3\)；此外，
\end{enumerate}

\[
\left(\operatorname{Aut} ^ {\prime} (\mathfrak {g}): \operatorname{Aut} _ {0} (\mathfrak {g})\right) = 2.
\]

元素 \(s \in \operatorname{Aut}(\mathfrak{g})\) 属于 \(\operatorname{Aut}'(\mathfrak{g})\)，当且仅当 \(\sigma \circ s\) 与 \(\sigma\) 等价（这是因为 \(\varpi_{1}\) 是 \(\sigma\) 的最高权）。如第2.VII节中那样，我们得到 \(\operatorname{Aut}'(\mathfrak{g})\) 可与 \(\Sigma/k^{*}\) 识别，其中 \(\Sigma\) 是关于 \(\Psi\) 的 V 的相似变换群。

令 \(s \in \Sigma\)，令 \(\lambda(s)\) 为 s 的乘子。如果 \(s\) 是正相似，则 \(\det(s) = \lambda(s)^l\)；如果 \(s\) 是逆相似，则 \(\det(s) = -\lambda(s)^l\)。因此，\(s\) 也是正相似；正相似构成 \(\Sigma_0\) 的一个指数为 2 的子群，并且 \(\Sigma\) 满足 \(\Sigma_0 \supset k^*\)。群 \(\Sigma_0 / k^*\) 等于 \(\operatorname{Aut}_0(\mathfrak{g})\)（它是 \(\operatorname{Aut}'(\mathfrak{g}) = \Sigma / k^*\) 的子群）。事实上，只需在 \(k\) 代数闭时验证这一点：此时 \(\operatorname{Aut}_0(\mathfrak{g}) = \operatorname{Aut}_e(\mathfrak{g})\) 等于其换位子群（...），所以包含于 \(\Sigma_0 / k^*\)；而二者在 \(\Sigma / k^*\) 中的指数都为 2，故二者相等。

另一方面，正如第3.VII节中那样，有一个正合序列

\[
1 \longrightarrow \mathbf {S O} (\varPsi) / \{1, - 1 \} \longrightarrow \varSigma_ {0} / k ^ {*} \longrightarrow k ^ {*} / k ^ {* 2} \longrightarrow 1.
\]

将 \(\mathbf{SO}(\varPsi)/\{1,-1\}\) 识别为 \(\Sigma_{0}/k^{*}=\mathrm{Aut}_{0}(\mathfrak{g})\) 的一个子群。由于 \(k^{*}/k^{*2}\) 是交换群，有 \(\mathrm{Aut}_{e}(\mathfrak{g})\subset\mathbf{SO}(\varPsi)/\{1,-1\}\)。事实上，可以证明（习题11），\(\mathrm{Aut}_{e}(\mathfrak{g})\) 等于 \(\mathbf{SO}(\varPsi)/\{1,-1\}\) 的既约正交群 \(\mathbf{O}_{0}^{+}(\varPsi)\) 在 \(\varPsi\) 中的像（《代数》第 IX 章，§9，第5节）。

\begin{enumerate}
\def\labelenumi{(\Roman{enumi})}
\setcounter{enumi}{7}
\tightlist
\item
典则双线性型 \(\varPhi_{\mathrm{R}}\) 在 \(\mathfrak{h}^*\) 上为
\end{enumerate}

\[
\Phi_ {\mathrm{R}} \left(\xi_ {1} \varepsilon_ {1} + \dots + \xi_ {l} \varepsilon_ {l}, \xi_ {1} ^ {\prime} \varepsilon_ {1} + \dots + \xi_ {l} ^ {\prime} \varepsilon_ {l}\right) = \frac {1}{4 (l - 1)} \left(\xi_ {1} \xi_ {1} ^ {\prime} + \dots + \xi_ {l} \xi_ {l} ^ {\prime}\right)
\]

（第 VI 章，§4，第8.V节）。因此，Killing 型在 \(\mathfrak{h}\) 上的限制为

\[
\Phi \left(\xi_ {1} H _ {1} + \dots + \xi_ {l} H _ {l}, \xi_ {1} ^ {\prime} H _ {1} + \dots + \xi_ {l} ^ {\prime} H _ {l}\right) = 4 (l - 1) \left(\xi_ {1} \xi_ {1} ^ {\prime} + \dots + \xi_ {l} \xi_ {l} ^ {\prime}\right).
\]

\begin{enumerate}
\def\labelenumi{(\Roman{enumi})}
\setcounter{enumi}{8}
\tightlist
\item
回顾由公式（11）定义的 \(X_{\alpha}\)（\(\alpha \in \mathbb{R}\)）。容易验证，对于 \([X_{\alpha}, X_{-\alpha}] = -H_{\alpha}\)，有 \(\alpha \in \mathbb{R}\)。另一方面，映射 \(\theta : a \mapsto -^{t}a\) 是 \(\mathfrak{g}\) 的自同构，且对所有 \(\theta(X_{\alpha}) = X_{-\alpha}\) 有 \(\alpha \in \mathbb{R}\)。因此，\((X_{\alpha})_{\alpha \in \mathbb{R}}\) 是 \((\mathfrak{g}, \mathfrak{h})\) 中的谢瓦莱系。
\end{enumerate}

假设 \(k = \mathbf{Q}\)。嘉当子代数 \(\mathfrak{h}\) 在 \(l\) 为奇数时有三个许可格，在 \(l\) 为偶数时有四个许可格。特别地，格 \(\mathcal{H}\) 由 \(H_i\) 生成，且是 \(\mathfrak{g}\) 的许可格。因此，\(\mathfrak{o}_S(2l,\mathbf{Z})\) 是谢瓦莱序 \(\mathcal{H} + \sum \mathbf{Z}.X_\alpha\)，属于 \(\mathfrak{g}\)。由于 \(X_{\alpha}^{2} = 0\) 对所有 \(\alpha \in \mathrm{R}\) 都有，可知格 \(\mathcal{V}\) 位于 \(\mathrm{V}\) 中，由 Witt 基 \((e_i)\) 生成；它是 \(\mathrm{V}\) 中关于 \(\mathfrak{o}_S(2l,\mathbf{Z})\) 的容许格。\(\bigwedge^r\mathcal{V}\) 在 \(\bigwedge^r\mathrm{V}\) 中也如此。

另一方面，若取 \(\mathrm{P}(\mathrm{R}^{\vee}) = \mathbf{Z}\) 作为许可格，并取 \(\frac{1}{2} \sum_{i=1}^{l} H_i + \mathcal{H}\) 作为谢瓦莱序，则可知 \(\mathcal{G} = \mathrm{P}(\mathrm{R}^{\vee}) + \sum \mathbf{Z}. X_{\alpha}\) 中的 \(\bigwedge^r \mathrm{V}\) 只有在 \(r\) 为偶数时才有容许格；此时 \(\bigwedge^r \mathcal{V}\) 是容许格。

考虑约化李代数 \(\mathfrak{o}(\varPsi) + \mathbf{Q}.1\)；立即可见，格 \(\tilde{\mathcal{G}} = (\mathfrak{o}(\varPsi) + \mathbf{Q}.1) \cap \mathfrak{gl}(2l, \mathbf{Z})\) 是一个谢瓦莱序。谢瓦莱序 \(\mathcal{G}\) 是 \(\tilde{\mathcal{G}}\) 平行于中心 \(\mathfrak{o}(\varPsi)\) 投影到 \(\mathbf{Q}.1\) 上所得的序。

最后，如第2节中那样可见，由 \(N_{+}\)（分别由 \(N_{-}\)）生成的格 \(e_{-i_{1}}\wedge\cdots\wedge e_{-i_{2k}}\)（分别 \(e_{-i_{1}}\wedge\cdots\wedge e_{-i_{2k+1}}\)）是谢瓦莱序 \(\mathrm{Q}(\mathrm{R}^{\vee})+\sum_{\alpha\in\mathbb{R}}\mathbf{Z}.X_{\alpha}\) 的半旋量表示的容许格。另一方面，\(N_{+}\) 和 \(N_{-}\) 对 \(\mathfrak{o}_{S}(2l,\mathbf{Z})\) 没有容许格。

\plainsection{表 1}

我们给每个基本权关联数 1（分别在相应的单表示为正交型时），以及 -1、0（分别在相应表示为辛型、既非正交型也非辛型时）。这个数的计算实质上已在§13中给出，适用于 \(\mathrm{A}_l, \mathrm{B}_l, \mathrm{C}_l, \mathrm{D}_l\) 型。对于 \(\mathrm{E}_6, \mathrm{E}_7, \mathrm{E}_8, \mathrm{F}_4, \mathrm{G}_2\) 型，结果也列于下方（只需应用§7，命题12，以及第 VI 章，§4，第9.VI、9.XI、10.VI、10.XI、11.VI、11.XI、12.VI、12.XI、13.VI、13.XI节）。

\[
\begin{array}{c c} \mathrm{A} _ {l} (l \geq 1) & \mathrm{B} _ {l} (l \geq 2) \\ \varpi_ {r} \left\{ \begin{array}{l l} 0 & r \neq \frac {l + 1}{2} \\ (- 1) ^ {r} & r = \frac {l + 1}{2} \end{array} \right. & \begin{array}{c c} \varpi_ {r} & 1 \quad r \neq l \\ \varpi_ {l} & (- 1) ^ {l (l + 1) / 2} \end{array} \end{array}
\]

\[
\begin{array}{c c} \varpi_ {r} & 1 \quad r \neq l - 1, l \\ \varpi_ {l} \text {and} \varpi_ {l - 1} & \left\{ \begin{array}{l l} 0 & \text {if l is odd} \\ (- 1) ^ {l / 2} & \text {if l is even} \end{array} \right. \end{array}
\]

\[
\begin{array}{c c c c c c c c c c} \mathrm{E} _ {6} & & \mathrm{E} _ {7} & & \mathrm{E} _ {8} & & \mathrm{F} _ {4} & & \mathrm{G} _ {2} \\ \varpi_ {1} & 0 & \varpi_ {1} & 1 & \varpi_ {1} & 1 & \varpi_ {1} & 1 & \varpi_ {1} & 1 \\ \varpi_ {2} & 1 & \varpi_ {2} & - 1 & \varpi_ {2} & 1 & \varpi_ {2} & 1 & \varpi_ {2} & 1 \\ \varpi_ {3} & 0 & \varpi_ {3} & 1 & \varpi_ {3} & 1 & \varpi_ {3} & 1 \\ \varpi_ {4} & 1 & \varpi_ {4} & 1 & \varpi_ {4} & 1 & \varpi_ {4} & 1 \\ \varpi_ {5} & 0 & \varpi_ {5} & - 1 & \varpi_ {5} & 1 \\ \varpi_ {6} & 0 & \varpi_ {6} & 1 & \varpi_ {6} & 1 \\ & & \varpi_ {7} & - 1 & \varpi_ {7} & 1 \\ & & & & \varpi_ {8} & 1 \end{array}
\]